% concatenated from Suspension-Bridge-Cattaneo-Heat.tex
\documentclass[11pt,reqno]{article}
\usepackage{diagbox}
\usepackage{amsfonts}
\usepackage{bbm}
\usepackage{bm}
\usepackage{pifont}
\usepackage{mathptmx}
\usepackage{amsmath}
\usepackage{mathrsfs}
\usepackage{hyperref}
\usepackage{amssymb}
\usepackage{amsmath,color}
\usepackage{amsthm}
\usepackage{amstext}
\usepackage[numbers,sort&compress]{natbib}
\usepackage{amsopn}
\usepackage{texdraw}
\usepackage{graphicx}
\usepackage{multirow}
\usepackage{lscape}
\usepackage{float}
\usepackage[palatino,gill,courier]{altfont}
\usepackage[left=3cm,right=3cm,top=3cm,bottom=3cm]{geometry}
\usepackage{esint}
\usepackage{setspace}

%\documentclass[10pt]{amsart}
%\usepackage{amsmath,amssymb,latexsym,esint,cite,mathrsfs}
%\usepackage{verbatim,wasysym}
%\usepackage[left=2.65cm,right=2.65cm,top=3.15cm,bottom=3.15cm]{geometry}

%\usepackage[colorlinks=true,urlcolor=blue, citecolor=red,linkcolor=blue,
%linktocpage,pdfpagelabels, bookmarksnumbered,bookmarksopen]{hyperref}

%\usepackage[hyperpageref]{backref}
\usepackage[english]{babel}
\usepackage{mathrsfs,amsmath}
\usepackage{color}
\usepackage{framed}
\allowdisplaybreaks



\usepackage{color} %make use of color features
\usepackage{enumerate} %utilizar listas con enumeraciones.
% estilo de enunciados de teoremas de la ams
\newtheorem{theorem}{Theorem}[section] %[section]: Enumera seg?n la secci?n
\newtheorem{proposition}[theorem]{Proposition}
\newtheorem{lemma}[theorem]{Lemma}
\newtheorem{example}{Example}[section]
\newtheorem{corollary}[theorem]{Corollary}
\newtheorem{definition}[theorem]{Definition}
\newtheorem{assumption}[theorem]{Assumption}
\newtheorem{remark}[theorem]{Remark}
\newtheorem{notation}[theorem]{Notation}
\def\theequation{\thesection.\arabic{equation}}
\newcommand{\R}{\mathbb{R}}
\newcommand{\C}{\mathbb{C}}
\newcommand{\M}{\mathcal{M}}
\newcommand{\dt}{\tilde}
\newcommand{\ds}{\displaystyle}
\newcommand{\be}{\beta}
\newcommand{\al}{\alpha}
\newcommand{\ta}{\theta}
\newcommand{\Hs}{H^1_*}
\newcommand{\Lsq}{L^2_*}
\newcommand{\D}{\mathcal{D}}
\newcommand{\diff}{\mathrm{d}}
\newcommand{\lk}{\left(}
\newcommand{\rk}{\right)}
\let\Section=\section
\def\section{\setcounter{equation}{0}\Section}
\oddsidemargin-0.3in \evensidemargin-0.3in\topmargin -0.3in \textwidth 7.2in \textheight 9.1in
\baselineskip=20pt
\parskip=2mm
\parindent=20pt
\setstretch{1.1}
\title{Well-posedness and Stability Analysis of Suspension Bridge Models Coupled with Cattaneo Heat Conduction: The Role of Viscoelastic Memory}
\author{Jun Zhou\footnote{Corresponding author. Email: jzhouwm@163.com} ~~~~~~ Jiamin He\footnote{Email: 1551937477@qq.com}\\\\{School of Mathematics and Statistics, Southwest University},\\
		{Chongqing 400715, People's Republic of China}}
\date{}
\date{}
\begin{document}
	\maketitle
	\begin{abstract}
This paper focuses on investigating a class of suspension bridge models coupled with the Cattaneo heat conduction law, with special attention paid to two distinct scenarios: systems with viscoelastic memory and those without. Operator semigroup theory is adopted as the core mathematical tool to conduct a comprehensive analysis of the models' dynamic behaviors. For the memoryless suspension bridge system ($\delta = 0$), we first establish its well-posedness (i.e., the existence and uniqueness of solutions). Further stability analysis reveals that this system exhibits polynomial decay behavior, where the decay rate depends on the structural parameter relationship of the bridge deck: a decay rate of $t^{-1/2}$ is achieved when $\frac{\rho_1}{\kappa} = \frac{\rho_2}{b}$, and the rate slows down to $t^{-1/4}$ when $\frac{\rho_1}{\kappa} \neq \frac{\rho_2}{b}$. Moreover, exponential stability is proven to be unachievable under specific coefficient configurations (i.e., $\chi_0 \neq 0$, or $\chi_0 = 0$ and $\chi_1 = 0$, with $\chi_0, \chi_1$ being parameter-dependent coefficients). For the counterpart system incorporating viscoelastic memory ($\delta = 1$), the well-posedness of solutions is also verified, and more importantly, the system is shown to achieve exponential decay. This indicates that viscoelastic memory significantly enhances system stability and accelerates energy dissipation compared to the memoryless case. By systematically exploring the regulatory role of viscoelastic memory in the stability of suspension bridge systems under the framework of Cattaneo thermal conduction, this study enriches and extends the existing research on  suspension bridge models with viscoelastic memory.
		\\\\
		\textbf{Keywords}:  Suspension bridge; Viscoelastic memory; Cattaneo heat conduction; Well-posedness; Stability\\\\
\textbf{MSC}:  35B35; 35B40; 35L70; 74D05; 80A20
	\end{abstract}
\section{Introduction}
As a bridge type boasting exceptional spanning capability, the grand structure of a suspension bridge embodies intricate mechanical behaviors. Consequently, developing accurate mathematical models is of utmost importance. In a typical suspension bridge model, the main cable is idealized as an elastic string to characterize its vertical vibrations, whereas the bridge deck is modeled based on the Timoshenko beam theory \cite{timoshenko1921lxvi}. The main cable and bridge deck are connected through a set of linear elastic springs. These springs simulate the suspenders, thereby forming a complex interactive system. To analyze the stability of this system, damping mechanisms (such as internal friction or thermoelastic effects) are frequently integrated into research studies.

In this paper, we consider the following suspension bridge system, which incorporates the Cattaneo heat conduction law as a key component
	\begin{align}\label{model1.1}
		\begin{cases}
			\rho u_{tt} - \alpha u_{xx} - \beta(\varphi - u) + \gamma u_t = 0 & \text{in } (0,L) \times \mathbb{R}^+, \\
			\rho_1 \varphi_{tt} - \kappa(\varphi_x + \psi)_x + \beta (\varphi - u) + m \theta_x = 0 & \text{in } (0,L) \times \mathbb{R}^+, \\
			\rho_2 \psi_{tt} - b \psi_{xx} + \kappa(\varphi_x + \psi) + \delta \int_0^\infty g(s) \psi_{xx}(t-s)  ds - m \theta = 0 & \text{in } (0,L) \times \mathbb{R}^+, \\
			\rho_3 \theta_t + q_x + m (\varphi_x + \psi)_t = 0 & \text{in } (0,L) \times \mathbb{R}^+, \\
			\tau q_t + \sigma q + \theta_x = 0& \text{in } (0,L) \times \mathbb{R}^+,
		\end{cases}
	\end{align}
where the parameter $\delta$ takes two distinct values to represent different viscoelastic conditions: $\delta = 0$ corresponds to the case \textit{without viscoelastic memory}; $\delta = 1$ corresponds to the case \textit{with viscoelastic memory}.

 This suspension bridge system is subject to the following homogeneous boundary conditions, which reflect the physical constraints at the two ends of the bridge (denoted by $x=0$ and $x=L$)
	\begin{align*}
		u_x(0,t) = u_x(L,t) = \varphi_x(0,t) = \varphi_x(L,t) = \psi(0,t) = \psi(L,t) = \theta(0,t) = \theta(L,t) = 0.
	\end{align*}

The key physical quantities in the aforementioned system are defined based on their functional roles in describing the dynamic (vibrational) and thermal behaviors of the suspension bridge: $u$ represents the vertical displacement of the main cable (captures the vertical vibrational response of the main cable under external or internal forces); $\varphi$ represents the transverse displacement of the bridge deck beam's cross-section (describes the lateral displacement of the centroid of the bridge deck beam's cross-section); $\psi$ represents the rotation angle of the bridge deck beam's cross-section (characterizes the rotational deformation of the bridge deck beam's cross-section during vibration); $q$ represents the heat flux (quantifies the rate of thermal energy transfer per unit area within the bridge system, a core variable in heat conduction analysis); $\theta$  represents the temperature deviation relative to the ambient temperature (reflects the difference between the system's local temperature and the surrounding environmental temperature).

Each parameter in the model is associated with specific material properties, structural stiffness, or physical effects of the suspension bridge components. They are classified as follows for clarity: $\rho$ denotes the density of the main cable (mass per unit volume of the main cable material, determining the cable's inertial response to vibration); $\alpha$  denotes the elastic modulus of the main cable (measure of the main cable's stiffness against elastic deformation; a higher value indicates greater resistance to stretching/compression); $\gamma$ denotes the internal damping of the main cable (accounts for energy dissipation within the main cable due to internal friction during vibration, preventing infinite amplitude growth); $\rho_1$ denotes the density of the bridge deck beam (mass per unit volume of the bridge deck beam material, influencing the beam's translational inertial behavior); $\kappa$ denotes the shear stiffness of the bridge deck beam (describes the bridge deck beam's ability to resist shear deformation, critical for capturing transverse shear effects in beam vibration); $\rho_2$ denotes the rotational moment of inertia of the bridge deck beam's cross-section (quantifies the beam's resistance to rotational acceleration, essential for modeling rotational vibration of the cross-section); $b$ denotes the bending stiffness of the bridge deck beam (characterizes the bridge deck beam's resistance to bending deformation; depends on the beam's material and cross-sectional geometry); $\beta$ denotes the stiffness of the suspenders (stiffness of the vertical suspenders connecting the main cable and the bridge deck beam; governs the force transfer between these two components); $m$ denotes the material-dependent coupling coefficient (quantifies the interactive effect between the structural deformation field and the thermal field, reflecting thermo-mechanical coupling); $\rho_3$ denotes the thermal inertia (represents the system's resistance to changes in temperature; a higher value means slower temperature variation); $\tau$ denotes the relaxation time (a key parameter in the Cattaneo heat conduction law, describing the delay in the response of heat flux to temperature gradients, distinguishing it from Fourier's law); $\sigma$ denotes the thermal conductivity (measure of the material's ability to conduct heat; a higher value indicates more efficient thermal energy transfer).

To ensure the physical rationality and mathematical well-posedness of the model, we impose the following assumptions on the parameters:  $\rho, \rho_1, \rho_2, \rho_3, \alpha, \kappa, b, \tau, \sigma, \beta, \gamma > 0$; the thermo-mechanical coupling coefficient $m$ is a non-zero real number, i.e., $m \in \mathbb{R} \setminus \{0\}$ (non-zero to reflect the existence of thermo-mechanical interaction). The specific definition of the memory kernel function $g(s)$, which describes the viscoelastic memory effect of the system when $\delta=1$, will be provided in subsequent sections, along with corresponding regularity conditions to guarantee the solvability of the integral term.
	
Raposo et al. \cite{MR4688212} investigated a suspension bridge model that incorporates only three types of internal damping: main cable damping, deck transverse damping, and deck rotational damping. For $(x,t) \in (0,L)\times \mathbb{R}^+$, the model is governed by the following system
	\begin{align*}
		\begin{cases}
			u_{tt} - \alpha u_{xx} - \beta(\varphi - u) + \gamma_1 u_t = 0,  \\
			\rho_1 \varphi_{tt} - \kappa (\varphi_x + \psi)_x + \beta (\varphi - u) + \gamma_2 \varphi_t = 0, \\
			\rho_2 \psi_{tt} - b \psi_{xx} + \kappa (\varphi_x + \psi) + \gamma_3 \psi_t = 0,\\
			u(0,t)=u(L,t)=\varphi(0,t)=\varphi(L,t)=\psi(0,t)=\psi(L,t)=0,
		\end{cases}
	\end{align*}
where they proved that the system achieves exponential stability relying solely on internal frictional damping.
	
Ivana Bochicchio et al. \cite{MR4149715} investigated a thermoelastically coupled suspension bridge model, where the main cable and deck are coupled via suspender stiffness. The model incorporates two types of internal damping: main cable damping and deck transverse damping and accounts for Fourier heat conduction effects. For $(x,t) \in (0,L)\times \mathbb{R}^+$, the model is formulated as the following system
	\begin{align*}
		\begin{cases}
			\varrho u_{tt} - \alpha u_{xx} - \beta(\varphi - u) + \gamma_1 u_t = 0, \\
			\varrho_1 \varphi_{tt} - \kappa(\varphi_x + \psi)_x + \beta(\varphi - u) + \gamma_2 \varphi_t = 0, \\
			\varrho_2 \psi_{tt} - b\psi_{xx} + \kappa(\varphi_x + \psi) - m\theta_x = 0, \\
			\theta_t - \tau \theta_{xx} - m\psi_{xt} = 0, \\
			u(0,t)=u(L,t)=\varphi(0,t)=\varphi(L,t)=\psi(0,t)=\psi_x(L,t)=\theta(0,t)=\theta(L,t)=0.
		\end{cases}
	\end{align*}
They proved the exponential stability of this thermoelastically coupled model: specifically, the total energy of the system decays at an exponential rate under the combined action of internal damping and thermal diffusion.

	
L. G. R. Miranda et al. \cite{MR4887867} investigated a thermoelastically coupled suspension bridge system that incorporates main cable internal damping and Fourier heat conduction effects, with the thermal dissipation mechanism acting on the shear force. For $(x,t) \in (0,L)\times \mathbb{R}^+$, the system is described by the following equations
	\begin{align}\label{eq1.3}
		\begin{cases}
			\rho u_{tt} - \alpha u_{xx} - \beta (\varphi - u) + \gamma u_t = 0,  \\
			\rho_1 \varphi_{tt} - \kappa (\varphi_x + \psi)_x + \beta(\varphi - u) + m\theta_x = 0 ,  \\
			\rho_2 \psi_{tt} - b \psi_{xx} + \kappa (\varphi_x + \psi) - m\theta = 0, \\
			\theta_t - \tau \theta_{xx} + m(\varphi_x + \psi)_t = 0,\\
			u_x(0,t)=u_x(L,t)=\varphi_x(0,t)=\varphi_x(L,t)=\psi(0,t)=\psi(L,t)=\theta(0,t)=\theta(L,t)=0.
		\end{cases}
	\end{align}
	They found that when $\frac{\kappa}{\rho_1} = \frac{b}{\rho_2}$, the system exhibits exponential stability; otherwise, it does not exhibit exponential stability and instead possesses polynomial stability with a decay rate of $t^{-\frac{1}{2}}$.

	
For type III thermoelasticity, Meriem Chabekh et al. \cite{MR4848548} analyzed the following suspension bridge model
	\begin{align*}
		\begin{cases}
			\rho u_{tt} - \alpha u_{xx} - \beta (\varphi - u) + \gamma_1 u_t = 0, \\
			\rho_1 \varphi_{tt} - \kappa (\varphi_x + \psi)_x + \beta (\varphi - u) + \gamma_2 \varphi_t + m \theta_x = 0,  \\
			-b \psi_{xx} + \kappa (\varphi_x + \psi) = 0, \\
			\rho_3 \theta_{tt} - \tau \theta_{xx} + m \varphi_{xtt} - a \theta_{xxt} = 0, \\
			u(0, t) = u(L, t) = \varphi (0, t) = \varphi (L, t) = \psi (0, t) = \psi (L, t) = \theta (0, t) = \theta (L, t) = 0,
		\end{cases}
	\end{align*}
where $(x,t) \in (0,L)\times \mathbb{R}^+$. Using the energy method, they proved that the system's energy decays exponentially over time.
	
	Soh Edwin Mukiawa et al. \cite{MR4539428} investigated a thermoelastic Timoshenko beam system with suspenders, weakly nonlinear internal damping (with time-varying coefficients), and time-varying delay. For $(x,t) \in (0,L) \times \mathbb{R}^+$, the system is formulated as follows
\begin{align*}
    \begin{cases}
        \rho u_{tt}- \alpha u_{xx} - \beta (\varphi - u)+ \gamma_1 a(t) g_1 (u_t(x,t)) + \gamma_2 a(t) g_2 (u_t(x,t - \tau(t))) = 0, \\
        \rho_1 \varphi_{tt} - \kappa(\varphi_x + \psi)_x + \beta (\varphi - u)+ \gamma_3 \varphi_t = 0, \\
        \rho_2 \psi_{tt} - b \psi_{xx}+ \kappa(\varphi_x + \psi) - m \theta_x = 0, \\
        \rho_3 \theta_t - \tau \theta_{xx} - m \psi_{tx}= 0, \\
        u(0,t)=u(L,t)=\varphi_{x}(0,t)=\varphi(L,t)=\psi(0,t)=\psi_{x}(L,t) = \theta_{x}(0,t)=\theta(1,t)=0.
    \end{cases}
\end{align*}
From their analysis, they derived a general stability result, which is expressed as
\[
    E(t) \leq \delta_1 \chi_1^{-1} \left( \delta_2 \int_0^t a(s) \, ds + \delta_3 \right), \quad \forall t \geq 0.\]
	
When viscoelastic memory effects are accounted for, Soh Edwin Mukiawa et al. \cite{MR4913141} investigated and proposed the following coupled suspension bridge system, which incorporates material memory effects and adheres to the Gurtin-Pipkins heat conduction law. For $(x,t) \in (0,L) \times \mathbb{R}^+$, the system is defined as
\begin{align*}
    \begin{cases}
        \rho u_{tt} + \alpha u_{xxxx} + \beta[u - \varphi]^+ - \int_0^{\infty} g(s) u_{xxxx}(x,t-s) \, ds = 0,  \\
        \rho_1 \varphi_{tt} - \kappa \varphi_{xx} - \beta [u - \varphi]^+ + m \theta_x = 0,  \\
        \rho_3 \theta_t - \delta \int_0^{\infty} h(s) \theta_{xx}(x,t-s) \, ds + m \varphi_{xt} = 0,  \\
        u(0,t) = u_{xx}(0,t) = \varphi_x(0,t) = \theta(0,t) =
        u(L,t) = u_{xx}(L,t) = \varphi(L,t) = \theta_x(L,t) = 0.
    \end{cases}
\end{align*}
From their analysis of the system, they derived an estimate for the system's energy, which is expressed as
\begin{align*}
    E(t) \leq
    \begin{cases}
        K e^{-\alpha t}, & \text{if } p = 1, \\
        K (1 + t)^{-\frac{1}{2p-2}}, & \text{if } 1 < p < \frac{3}{2}.
    \end{cases}
\end{align*}
	
	Mohammad M. Al-Gharabli et al. \cite{MR4891851} investigated a one-dimensional coupled suspension bridge system model that incorporates both viscoelastic damping and frictional damping with the latter modulated by a time-dependent coefficient. Specifically, the model is expressed as follows
\begin{align*}
    \begin{cases}
        u_{tt} + u_{xxxx} + [u - \varphi]^+ + \beta(t)\phi(u_t) = 0, \\
        \varphi_{tt} - \varphi_{xx} - [u - \varphi]^+ + \int_0^t g(t-s)\varphi_{xx}(s)ds = 0,  \\
        u(0, t) = u_{xx}(0, t) = u(L, t) = u_{xx}(L, t) = \varphi(0, t) = \varphi(L, t) = 0,
    \end{cases}
\end{align*}
where the spatial-temporal domain is defined as $(x,t) \in (0,L)\times (0,T)$ ($T$ denotes the finite time horizon for analysis). From their investigation, they derived an estimate for the system's energy, which is given by
\begin{align*}
    E(t) \leq c_1 \Phi_1^{-1} \left( m_2 \int_0^t \xi(s)  ds + m_3 \right), \quad \forall t \geq 0.
\end{align*}
	
	 Mounir Afilal et al. \cite{MR4342054} investigated a one-dimensional coupled suspension bridge system model with viscoelastic memory, which is formulated as follows
\begin{align*}
    \begin{cases}
        u_{tt} + u_{xxxx} - \int_0^t g_1 (t-s) u_{xxxx}(s)  ds + \left( \int_0^L (u^2 + \varphi^2)  dx \right) u + 2 \left( \int_0^L u \varphi  dx \right)\varphi = 0,\\
        \varphi_{tt} - \varphi_{xx} + \int_0^t g_2 (t-s) \varphi_{xx}(s)  ds + 2 \left( \int_0^L u \varphi  dx \right) u + \left( \int_0^L (u^2 + \varphi^2)  dx \right) \varphi= 0, \\
        u(0,t)=u(L,t)=u_{xx}(0,t)=u_{xx}(L,t)=\varphi(0,t)=\varphi(L,t)=0,
    \end{cases}
\end{align*}
where $(x,t) \in (0,L)\times \mathbb{R}^+$. Using the energy method, they derived a decay estimate for the system's energy, given by
\begin{align*}
    E(t) \leq C_0 H_4^{-1} \left( C_1 \int_{t_0}^{t} \xi(s)  ds \right), \quad \forall t \geq t_0.
\end{align*}

Furthermore, for the stability analysis of two-dimensional suspension bridges, several key studies have been conducted: for instance, Soh Edwin Mukiawa \cite{MR3754568} analyzed a two-dimensional fourth-order viscoelastic suspension bridge system and derived an energy estimate for it. Notably, Zayd Hajjej \cite{MR4252268} investigated a viscoelastic suspension bridge model incorporating nonlinear damping and a source term, also providing an energy estimate for the system. Building on this work, Adel M. Al-Mahdi et al. \cite{MR4694563} extended the results presented in \cite{MR4252268}.

Additionally, additional studies on the asymptotic behavior of suspension bridges are categorized below for reference: For models without memory terms, readers may refer to \cite{MR4875485, MR3023261, MR2142424, MR2514725}; for models that incorporate memory effects, the recommended references are \cite{MR4886591, MR3575770, MR3240285, MR3458828, MR3609015}.

The decision to focus on the aforementioned suspension bridge system \eqref{model1.1} stems from three key considerations, which address critical gaps identified in the existing literature reviewed above. First, regarding heat conduction modeling: most prior thermoelastically coupled models (e.g., Bochicchio et al. \cite{MR4149715}, Miranda et al. \cite{MR4887867}) adopt Fourier's law, which assumes instantaneous thermal propagation-an idealization that deviates from engineering reality, where heat conduction inherently exhibits time delay (especially in materials with low thermal conductivity, such as concrete used in bridge decks). By incorporating the Cattaneo heat conduction law (characterized by the relaxation time $\tau$ and heat flux $q$), our model accurately captures this thermal delay, making it more representative of the actual thermal behavior of suspension bridges. Second, in terms of viscoelastic memory representation: existing models either exclude memory effects entirely (Raposo et al. \cite{MR4688212}, Chabekh et al. \cite{MR4848548}) or fix memory as a permanent feature (Mukiawa et al. \cite{MR4913141}, Al-Gharabli et al. \cite{MR4891851}), while practical suspension bridges may experience varying viscoelastic behavior (e.g., temporary memory effects due to temperature fluctuations or long-term memory from material creep). Our model's adjustable parameter $\delta$ ($\delta=0$ for no memory, $\delta=1$ for memory) enables a unified analysis of both scenarios, filling the gap of \lq\lq switchable viscoelasticity\rq\rq in existing research. Third, from the perspective of multi-physical coupling comprehensiveness: the model integrates four core mechanisms-main cable vertical vibration (governed by $u$), bridge deck transverse/rotational vibration (based on Timoshenko beam theory, described by $\varphi$ and $\psi$), Cattaneo-type heat conduction (via $\theta$ and $q$), and thermo-mechanical coupling (through $m$), while also including main cable internal damping ($\gamma$) and suspender stiffness ($\beta$). This holistic framework avoids the simplifications of prior models (e.g., neglecting deck rotation, using linear heat conduction, or omitting damping), ensuring it can capture the intricate interactions between mechanical and thermal fields in real suspension bridges. In summary, the model \eqref{model1.1} not only addresses the limitations of existing studies but also provides a more accurate and flexible theoretical tool for analyzing the stability and dynamic behavior of suspension bridges, with direct implications for engineering design and safety assessment.

The main results of this paper concerning the asymptotic stability of the suspension bridge system \eqref{model1.1} are summarized as follows
\begin{enumerate}
  \item When $\delta = 0$ (i.e., the system without viscoelastic memory), the solution to problem \eqref{model1.1} exhibits polynomial decay behavior. More precisely, the decay rate depends on the relationship between the structural parameters of the bridge deck: if $\frac{\rho_1}{\kappa} = \frac{\rho_2}{b}$, the solution decays polynomially with a rate of $t^{-1/2}$; otherwise, when $\frac{\rho_1}{\kappa} \neq \frac{\rho_2}{b}$, the decay rate slows down to $t^{-1/4}$. Moreover, if $\chi_0 \neq 0$, or if $\chi_0 = 0$ and $\chi_1 = 0$, the solution to problem \eqref{model1.1} is not exponentially stable, where $\chi_0 = \kappa\left(\dfrac{\rho_2}{b} -  \dfrac{\rho_1}{\kappa}\right)\left(\rho_3 - \dfrac{\rho_2}{b\tau}\right) + \dfrac{m^2 \rho_2}{b}$ and $\chi_1 = -m^2 - \kappa\rho_3 + \dfrac{\kappa\rho_2}{\tau b}$.
  \item When $\delta = 1$ (i.e., the system with viscoelastic memory), the solution to problem \eqref{model1.1} achieves exponential decay. This indicates that the introduction of viscoelastic memory significantly enhances the stability of the system, leading to a much faster energy dissipation compared to the memoryless case.
\end{enumerate}
      The main results reveal a crucial insight into the role of viscoelastic memory in regulating the stability of thermoelastic suspension bridge systems. Firstly, the polynomial stability in the memoryless case ($\delta = 0$) is consistent with the findings of Miranda et al. \cite{MR4887867}, which confirms the rationality of our model in capturing the intrinsic dynamic characteristics of suspension bridges without viscoelastic effects. However, our result further refines the decay rate by identifying the critical parameter ratio $\frac{\rho_1}{\kappa} = \frac{\rho_2}{b}$ that governs the polynomial decay speed, which is attributed to the more accurate Cattaneo heat conduction law adopted in our model. Secondly, the exponential stability in the memory case ($\delta = 1$) demonstrates the dominant role of viscoelastic memory in promoting energy dissipation. This is because the memory kernel $g(s)$ introduces additional damping mechanisms that continuously dissipate the system's energy accumulated from historical deformations, thereby overcoming the limited dissipation capacity of thermoelastic effects alone. From an engineering perspective, these results provide valuable guidance for the design and maintenance of suspension bridges: for bridges in environments where viscoelastic memory effects are significant (e.g., long-span bridges with polymer materials), the inherent viscoelasticity can be leveraged to enhance structural stability; for memoryless structures, the parameter ratio $\frac{\rho_1}{\kappa}$ and $\frac{\rho_2}{b}$ should be carefully optimized to ensure an acceptable decay rate. Finally, it is worth noting that the exponential decay result does not require the parameter ratio condition, implying that viscoelastic memory can uniformly improve the stability of the system regardless of the initial structural parameter configuration.

	
The remainder of this paper is organized as follows. In Section \ref{sec2}, we focus on the case \lq\lq without\rq\rq the viscoelastic memory term, corresponding to $\delta = 0$. Specifically, Subsection \ref{sec2.1} is dedicated to proving the well-posedness of the system under this condition. In Subsection \ref{sec2.2}, we establish the polynomial stability of the system, followed by Subsection \ref{sec2.3}, where we further demonstrate that the system fails to exhibit exponential stability under certain specified conditions. Subsequently, Section \ref{sec3} is devoted to the analysis of the case \lq\lq with\rq\rq the viscoelastic memory term, i.e., $\delta = 1$. In this section, Subsection \ref{sec3.1} verifies the well-posedness of the memory-dependent system, and Subsection \ref{sec3.2} rigorously proves that the system achieves exponential stability in this scenario.
\section{Model Formulation without the Boltzmann Memory Term ($\delta = 0$)} \label{sec2}
When $\delta = 0$, the suspension bridge model \eqref{model1.1} reduces to the following memoryless form
\begin{align}\label{eq:without memory}
	\begin{cases}
		\rho u_{tt} - \alpha u_{xx} - \beta(\varphi - u) + \gamma u_t = 0 & \text{in } (0,L) \times \mathbb{R}^+, \\
		\rho_1 \varphi_{tt} - \kappa(\varphi_x + \psi)_x + \beta (\varphi - u) + m \theta_x = 0 & \text{in } (0,L) \times \mathbb{R}^+, \\
		\rho_2 \psi_{tt} - b \psi_{xx} + \kappa(\varphi_x + \psi) - m \theta = 0 & \text{in } (0,L) \times \mathbb{R}^+, \\
		\rho_3 \theta_t + q_x + m (\varphi_x + \psi)_t = 0 & \text{in } (0,L) \times \mathbb{R}^+, \\
		\tau q_t + \sigma q + \theta_x = 0& \text{in } (0,L) \times \mathbb{R}^+,
	\end{cases}
\end{align}
This memoryless system is subject to the following initial and boundary conditions: for $t \in \mathbb{R}^+$, the boundary conditions are
\begin{align}\label{without memory:boundary condition}
        u_x(0,t) = u_x(L,t) = \varphi_x(0,t) = \varphi_x(L,t) = \psi(0,t) = \psi(L,t) = \theta(0,t) = \theta(L,t) = 0,
\end{align}
and for $x \in (0, L)$, the initial conditions are
\begin{equation}\label{without memory:initial condition}
 \begin{cases}
  u(x,0) = u_0(x), \quad u_t(x,0) = v_0(x), \\
  \varphi(x,0) = \varphi_0(x), \quad \varphi_t(x,0) = \Phi_0(x), \\
  \psi(x,0) = \psi_0(x), \quad \psi_t(x,0) = \Psi_0(x), \\
  \theta(x,0) = \theta_0(x), \quad q(x,0) = q_0(x).
 \end{cases}
\end{equation}

We consider the standard Lebesgue and Sobolev spaces over the interval \((0,L)\), denoted by
\[
L^2 = L^2(0,L), \quad H^1 = H^1(0,L), \quad H_0^1 = H_0^1(0,L).
\]
The standard inner product and norm on \(L^2\) are denoted by \((\cdot,\cdot)\) and \(\|\cdot\|\), respectively.
We introduce the following subspaces
\begin{align*}
	L_*^2 = \Bigl\{ \omega \in L^2 : \int_0^L \omega(x)\,dx = 0 \Bigr\}, ~~H_*^1 := H^1 \cap L_*^2, ~~L_g^2(\mathbb{R}^+; H_0^1) := \Bigl\{ \omega : \sqrt{g}\,\omega \in L^2\bigl(\mathbb{R}^+; H_0^1\bigr) \Bigr\},
\end{align*}
where the last space is equivalently characterized by the finiteness of the integral
\[
\int_0^{\infty} g(s)\, \|\omega_x(s)\|^2 \, ds < \infty .
\]

By virtue of Poincar\'e's inequality, which asserts that for any \(\phi \in H_0^1\) or \(\phi \in H_*^1\),
\[
\|\phi\|^2 \le c_p \|\phi_x\|^2
\]
holds with a positive constant \(c_p > 0\), the following Hilbert space structures are induced
\begin{itemize}
	\item  The spaces \(H_0^1\) and \(H_*^1\) are Hilbert spaces when equipped with the inner product \((\cdot_x, \cdot_x)\) (i.e., the inner product induced by the derivative) and the corresponding norm \(\|\cdot_x\|\).
	
	\item The subspace \(L_*^2\) is a Hilbert space under the standard \(L^2\)-inner product \((\cdot,\cdot)\) and its associated norm \(\|\cdot\|\).
	
	\item  The space \(L_g^2(\mathbb{R}^+; H_0^1)\) is a Hilbert space with respect to the inner product
        \[
        (\cdot, \cdot)_{L_g^2(\mathbb{R}^+; H_0^1)} := \int_0^{\infty} g(s)\, (\cdot_x, \cdot_x) \, ds
        \]
        and the corresponding norm
        \[
        \|\cdot\|_{L_g^2(\mathbb{R}^+; H_0^1)} := \Bigl( \int_0^{\infty} g(s)\, \|\cdot_x\|^2 \, ds \Bigr)^{1/2}.
        \]
\end{itemize}

Let \(v = u_t\), \(\Phi = \varphi_t\), \(\Psi = \psi_t\). To reformulate the problem as an abstract Cauchy problem, we introduce the phase space
\[
\mathcal{H}_1 := H_*^1 \times L_*^2 \times H_*^1 \times L_*^2 \times H_0^1 \times L^2 \times L^2 \times L^2,
\]
which is a Hilbert space endowed with the inner product
\begin{align*}
	(U_1, U_2)_{\mathcal{H}_1} := &\rho(v_1, v_2) + \alpha (u_{1x}, u_{2x}) + \rho_1(\Phi_1, \Phi_2) + \rho_2(\Psi_1, \Psi_2) + b (\psi_{1x},\psi_{2x}) \\
	&+ \beta (\varphi_1 - u_1, \varphi_2 - u_2) + \kappa(\varphi_{1x} + \psi_1, \varphi_{2x} + \psi_2) + \rho_3(\theta_1, \theta_2) + \tau(q_1, q_2),
\end{align*}
and the associated norm
\[
\|U\|_{\mathcal{H}_1}^2 := \rho\|v\|^2 + \alpha\|u_x\|^2 + \rho_1\|\Phi\|^2 + \rho_2\|\Psi\|^2 + b\|\psi_x\|^2 + \beta\|\varphi - u\|^2 + \kappa\|\varphi_x + \psi\|^2 + \rho_3\|\theta\|^2 + \tau \|q\| ^2,
\]
where \(U := (u, v, \varphi, \Phi, \psi, \Psi, \theta, q) \in \mathcal{H}_1\) and \(U_i := (u_i, v_i, \varphi_i, \Phi_i, \psi_i, \Psi_i, \theta_i, q_i) \in \mathcal{H}_1\) for \(i = 1, 2\).

The system \eqref{eq:without memory}-\eqref{without memory:initial condition} can be rewritten as the following abstract Cauchy problem
\begin{equation}\label{model-abstract-1}
	\begin{cases}
		\dfrac{d}{dt}U = \mathcal{A}_1U, & t > 0, \\
		U(0) =  U_0 :=(u_0, v_0, \varphi_0, \Phi_0, \psi_0, \Psi_0, \theta_0, q_0),
	\end{cases}
\end{equation}
where \(\mathcal{A}_1: \mathcal{D}(\mathcal{A}_1) \subset \mathcal{H}_1 \to \mathcal{H}_1\) is a linear operator defined by \begin{equation}\label{definition:A-1}
	\mathcal{A}_1 U := \begin{pmatrix}
		v \\
		\rho^{-1} \left[ \alpha u_{xx} + \beta (\varphi - u) - \gamma v \right] \\
		\Phi \\
		\rho_1^{-1} \left[ \kappa (\varphi_x + \psi)_x - \beta (\varphi - u) - m \theta_x \right] \\
		\Psi \\
		\rho_2^{-1} \left[ b \psi_{xx} - \kappa (\varphi_x + \psi) + m\theta \right] \\
		\rho_3^{-1} \left[ -q_x - m (\Phi_x + \Psi) \right] \\
		\tau^{-1}(-\sigma q - \theta_x)
	\end{pmatrix}
\end{equation}
with domain
\begin{align*}
	\mathcal{D}(\mathcal{A}_1) := \bigg\{ U \in \mathcal{H}_1 :
	v, \Phi \in H_*^1;  u_x, \varphi_x, \theta, \Psi \in H_0^1;
	\psi \in H^2; q \in H^1
	\bigg\}.
\end{align*}

The main results of this section are as follows
\begin{theorem}[Well-posedness]\label{Well-Posedness-1}
	The operator $\mathcal{A}_1$ defined in \eqref{definition:A-1} acts as the infinitesimal generator of a $C_0$-contraction semigroup $\{e^{t\mathcal{A}_1}\}_{t \ge 0}$. Accordingly, if $U_0 \in \mathcal{H}_1$, then Problem \eqref{model-abstract-1} has a unique mild solution $U(t) = e^{t\mathcal{A}_1}U_0$ belonging to $C(\mathbb{R}^+, \mathcal{H}_1)$; if $U_0 \in \mathcal{D} (\mathcal{A}_1)$ (the domain of $\mathcal{A}_1$), then $U(t) = e^{t\mathcal{A}_1}U_0$ is a classical solution to Problem \eqref{model-abstract-1} and lies in $C(\mathbb{R}^+, \mathcal{D}(\mathcal{A}_1)) \cap C^1(\mathbb{R}^+, \mathcal{H}_1)$.
\end{theorem}
\begin{theorem}[Polynomial stability]\label{Polynomial stability}
   Let \(\{e^{\mathcal{A}_1 t}\}_{t \geq 0}\) be the \(C_0\)-semigroup generated by the operator \(\mathcal{A}_1: \mathcal{D}(\mathcal{A}_1) \subset \mathcal{H}_1 \to \mathcal{H}_1\) (see \eqref{definition:A-1} for the operator definition and \(\mathcal{D}(\mathcal{A}_1)\) for its domain). Then \(\{e^{\mathcal{A}_1 t}\}_{t \geq 0}\) is polynomially stable, and the decay rate depends on the relationship between the structural parameters of the bridge deck
    \begin{enumerate}[(i)]
    	\item If \(\frac{\rho_1}{\kappa} = \frac{\rho_2}{b}\), there exists a positive constant \(C > 0\) (independent of \(t\)) such that
        \[
        \|e^{\mathcal{A}_1 t} U_0\|_{\mathcal{H}_1} \leq \frac{C}{t^{1/2}} \|U_0\|_{\mathcal{D}(\mathcal{A}_1)} \quad \text{for all } t > 0 \text{ and } U_0 \in \mathcal{D}(\mathcal{A}_1);
        \]
        \item If \(\frac{\rho_1}{\kappa} \neq \frac{\rho_2}{b}\), there exists a positive constant \(C > 0\) (independent of \(t\)) such that
        \[
        \|e^{\mathcal{A}_1 t} U_0\|_{\mathcal{H}_1} \leq \frac{C}{t^{1/4}} \|U_0\|_{\mathcal{D}(\mathcal{A}_1)} \quad \text{for all } t > 0 \text{ and } U_0 \in \mathcal{D}(\mathcal{A}_1),
        \]
    \end{enumerate}
    where $U_0 \in \mathcal{D}(\mathcal{A}_1):=\|\mathcal{A}_1U_0\|_{\mathcal{H}_1}$.
  \end{theorem}
  \begin{theorem}[Non-exponential stability]\label{Non-exponential Stability}
 Let the parameter-dependent coefficients $\chi_0$ and $\chi_1$ be defined as
        \[
        \chi_0 = \kappa\left(\dfrac{\rho_2}{b} - \dfrac{\rho_1}{\kappa}\right)\left(\rho_3 - \dfrac{\rho_2}{b\tau}\right) + \dfrac{m^2 \rho_2}{b}, \quad \chi_1 = -m^2 - \kappa\rho_3 + \dfrac{\kappa\rho_2}{\tau b}.
        \]
        If either $\chi_0 \neq 0$, or $\chi_0 = 0$ and $\chi_1 = 0$, then the $C_0$-semigroup $\{e^{\mathcal{A}_1 t}\}_{t \geq 0}$ generated by $\mathcal{A}_1$ is not exponentially stable.
  \end{theorem}
    \begin{remark}
  	It should be noted that for the case where $\chi_0 = 0$ and $\chi_1 \neq 0$, the exponential stability of the semigroup $\{e^{\mathcal{A}_1 t}\}_{t \geq 0}$ remains an open problem and requires further investigation.
  \end{remark}

\subsection{Well-Posedness}\label{sec2.1}
The proof of Theorem \ref{Well-Posedness-1} proceeds via an application of the Lumer-Phillips Theorem \cite{MR710486}. Consequently, the demonstration hinges on verifying the following two conditions of the theorem: (1) $\mathcal{A}_1$ is dissipative; (2) $\mathcal{A}_1$ is maximal. We now proceed by verifying these two conditions sequentially in Lemmas \ref{dissipative-1} and \ref{maximal-1}.

\begin{lemma}\label{dissipative-1}
	The operator $\mathcal{A}_1$ defined in \eqref{definition:A-1} is dissipative.
\end{lemma}
\begin{proof}
	For any $U \in \mathcal{D}(\mathcal{A}_1)$, by direct calculation and simplification via integration by parts, we obtain
	\begin{align*}
		{\rm Re}(\mathcal{A}_1U,U)_{\mathcal{H}_1} &= {\rm Re}\bigg\{
		\rho\left(\rho^{-1} \left[ \alpha u_{xx} + \beta (\varphi - u) - \gamma v \right], v\right) + \alpha\left(v_x, u_x\right)  + \rho_1\left(\rho_1^{-1} \left[ \kappa (\varphi_x + \psi)_x - \beta (\varphi - u) - m \theta_x \right], \Phi\right) \\
		&\quad + \rho_2\left(\rho_2^{-1} \left[ b \psi_{xx} - \kappa (\varphi_x + \psi) + m\theta \right], \Psi\right) + b\left(\Psi_x,\psi_x\right) + \beta\left(\Phi - v, \varphi - u\right)  \\
		&\quad + \kappa\left(\Phi_x + \Psi, \varphi_x + \psi\right)  + \rho_3\left(\rho_3^{-1} \left[ -q_x - m (\Phi_x + \Psi) \right], \theta\right) + \tau\left(\tau^{-1}(-\sigma q - \theta_x),q\right) \bigg\} \\
		&= -\gamma\|v\|^2  - \sigma \|q\|^2\\
&\le 0.
		\end{align*}
	Thus, $\mathcal{A}_1$ satisfies the dissipativity condition.
\end{proof}
\begin{lemma}\label{maximal-1}
	$0 \in \varrho(\mathcal{A}_1)$ (the resolvent set of $\mathcal{A}_1$). Consequently, $\mathcal{A}_1$ is maximal.
\end{lemma}
\begin{proof}
	First, we will prove for any $F = (f_1,f_2,f_3,f_4,f_5,f_6,f_7,f_8) \in \mathcal{H}_1$, there exists a unique $$U = (u,v,\varphi,\Phi,\psi,\Psi,\theta,q)\in \mathcal{D}(\mathcal{A}_1)$$ such that \begin{align*}
		\mathcal{A}_1U = F.
	\end{align*}
	From the definition of $\mathcal{A}_1$ \eqref{definition:A-1}, the expression can be reformulated as the following equations
	\begin{equation}\label{eq:resolvent-system-A-1}
		\begin{cases}
			v = f_1, \\
			\alpha u_{xx} + \beta (\varphi - u) - \gamma v = \rho f_2, \\
			\Phi = f_3, \\
			\kappa (\varphi_x + \psi)_x - \beta (\varphi - u) - m \theta_x = \rho_1 f_4, \\
			\Psi = f_5, \\
			b \psi_{xx} - \kappa (\varphi_x + \psi) + m\theta = \rho_2 f_6, \\
			-q_x - m (\Phi_x + \Psi) = \rho_3 f_7, \\
			-\sigma q - \theta_x = \tau f_8.
		\end{cases}
	\end{equation}
	From equations \eqref{eq:resolvent-system-A-1}$_{1,3,5}$, it follows that
	\begin{equation}\label{df1}
		v = f_1 \in H_*^1, \quad \Phi = f_3 \in H_*^1, \quad \Psi = f_5 \in H_0^1.
	\end{equation}
	Equation \eqref{eq:resolvent-system-A-1}$_8$ yields
	$q_x=-\frac{1}{\sigma}(\theta_{xx}+\tau f_{8x}).$ By substituting the above equation into equation \eqref{eq:resolvent-system-A-1}$_7$, we obtain
	\begin{align*}
		- \theta_{xx} = - \sigma m(\Phi_x + \Psi) - \sigma \rho_3 f_7 + \tau f_{8x}.
	\end{align*}
	Then combined with equations \eqref{eq:resolvent-system-A-1}$_{2,4,6}$, we derive
	\begin{align}\label{eq:resolvent-system-1-A-1}
		\begin{cases}
			-\alpha u_{xx} - \beta(\varphi - u) = -\rho f_2 - \gamma v=-\rho f_2 - \gamma f_1 := g_1 \in L_*^2, \\
			-\kappa(\varphi_x + \psi)_x + \beta(\varphi - u) + m\theta_x = -\rho_1 f_4 := g_2 \in L_*^2, \\
			- b \psi_{xx} + \kappa(\varphi_x + \psi) - m\theta = -\rho_2 f_6  := g_3 \in L^2, \\
			- \theta_{xx} = - \sigma m(\Phi_x + \Psi) - \sigma \rho_3 f_7 + \tau f_{8x}=- \sigma m(f_{3x} + f_5) - \sigma \rho_3 f_7 + \tau f_{8x}:=g_4 \in H^{-1}.
		\end{cases}
	\end{align}
	The weak formulation of equations \eqref{eq:resolvent-system-1-A-1} is defined as follows
	\begin{align}\label{eq:resolvent-system-2-A-1}
		\begin{cases}
			\alpha(u_x, \tilde{u}_x) - \beta(\varphi - u, \tilde{u}) = (g_1, \tilde{u}), \\
			\kappa(\varphi_x + \psi, \dt \varphi_x) + \beta(\varphi - u, \dt \varphi) - m(\theta, \dt \varphi_x) = (g_2, \dt \varphi), \\
			b(\psi_x, \dt \psi_x) + \kappa(\varphi_x + \psi, \dt \psi) - m(\theta, \dt \psi) = (g_3, \dt \psi), \\
			\left(\theta_x, \frac{m^2 c_p}{\kappa } \dt \theta_x\right) = \left(g_4, \frac{m^2 c_p}{\kappa } \dt \theta\right),
		\end{cases}
	\end{align}
	for any $(\tilde{u}, \tilde{\varphi}, \tilde{\psi}, \tilde{\theta}) \in V := H_*^1 \times H_*^1 \times H_0^1 \times H_0^1$. Summing each equation in \eqref{eq:resolvent-system-2-A-1} yields the bilinear form $\mathcal{L}: V \times V \to \mathbb{R}$ and the linear form $\mathcal{F}:V \to \mathbb{R}$, satisfying
	\begin{equation}\label{LF}
		\mathcal{L}\left(
		\begin{pmatrix}u\\ \varphi \\ \psi \\ \theta\end{pmatrix},
		\begin{pmatrix}\tilde{u}\\ \tilde{\varphi} \\ \tilde{\psi} \\ \tilde{\theta}\end{pmatrix}
		\right) = \mathcal{F}\left(
		\begin{pmatrix}\tilde{u}\\ \tilde{\varphi} \\ \tilde{\psi} \\ \tilde{\theta}\end{pmatrix}
		\right).
	\end{equation}
	where $\mathcal{L}$ and $\mathcal{F}$ are defined as
	\begin{equation}\label{L}
		\begin{split}
			\mathcal{L}\left(
			\begin{pmatrix}u\\ \varphi \\ \psi \\ \theta\end{pmatrix},
			\begin{pmatrix}\tilde{u}\\ \tilde{\varphi} \\ \tilde{\psi} \\ \tilde{\theta}\end{pmatrix}
			\right) &= \alpha (u_x, \tilde{u}_x) + \beta (\varphi - u, \tilde{\varphi} - \tilde{u}) + \kappa (\varphi_x + \psi, \tilde{\varphi_x} + \tilde{\psi}) \\
			&\quad - m(\theta, \tilde{\varphi_x}+ \tilde{\psi}) + b (\psi_x, \tilde{\psi_x})   + \frac{m^2 c_p}{\kappa} (\theta_x, \tilde{\theta_x}),
		\end{split}
	\end{equation}
	\begin{align}\label{F}
		\mathcal{F}\left(
		\begin{pmatrix}\tilde{u}\\ \tilde{\varphi} \\ \tilde{\psi} \\ \tilde{\theta}\end{pmatrix}
		\right) &= \left(g_1, \tilde{u}\right) + \left(g_2, \tilde{\varphi}\right) + \left(g_3, \tilde{\psi}\right) + \left(g_4, \frac{m^2 c_p}{\kappa }\tilde{\theta}\right) \notag \\
		&=  - \gamma (f_1, \tilde{u}) - \rho (f_2, \tilde{u}) - \rho_1 (f_4, \tilde{\varphi}) - \rho_2 (f_6, \tilde{\psi})  \notag \\
		&\quad  - \frac{\sigma m^3 c_p }{\kappa } (f_{3x} + f_5, \tilde{\theta}) - \frac{\sigma \rho_3 m^2 c_p }{\kappa } (f_7, \tilde{\theta}) - \frac{\tau m^2 c_p}{\kappa}(f_8, \tilde{\theta_x}).
	\end{align}
	Clearly, $\mathcal{L}$ and $\mathcal{F}$ are bounded. We now prove that $\mathcal{L}$ is coercive. In fact, for any $(u, \varphi, \psi, \theta) \in V$, by Young's inequality and Poincar\'e's inequality, we get
	\begin{align*}\label{L-1}
		\mathcal{L}\left(
		\begin{pmatrix}u\\ \varphi \\ \psi \\ \theta\end{pmatrix},
		\begin{pmatrix}u\\ \varphi \\ \psi \\ \theta\end{pmatrix}
		\right) &= \alpha \| u_x \|^2 + \beta \| \varphi - u \|^2 + \kappa \| \varphi_x + \psi \|^2 + b \| \psi_x \|^2 + \frac{m^2 c_p}{\kappa} \| \theta_x \|^2 - m (\theta, \varphi_x + \psi)\\
		&\geq \alpha \| u_x \|^2 + \beta \| \varphi - u \|^2 + \kappa \| \varphi_x + \psi \|^2 + b \| \psi_x \|^2 + \frac{m^2 c_p}{\kappa} \| \theta_x \|^2 - |m|\|\theta\|\|\varphi_x+\psi\|\\
		&\geq \alpha \| u_x \|^2 + \beta \| \varphi - u \|^2 + \kappa \| \varphi_x + \psi \|^2 + b \| \psi_x \|^2 + \frac{m^2 c_p}{\kappa} \| \theta_x \|^2 - \frac{\kappa}{2}\|\varphi_x+\psi\|^2 - \frac{m^2}{2\kappa}\|\theta\|^2\\
		&\geq \alpha \| u_x \|^2 + \beta \| \varphi - u \|^2 + \kappa \| \varphi_x + \psi \|^2 + b \| \psi_x \|^2 + \frac{m^2 c_p}{\kappa} \| \theta_x \|^2  - \frac{\kappa}{2}\|\varphi_x+\psi\|^2 - \frac{m^2 c_p}{2\kappa}\|\theta_x\|^2\\	
		&= \alpha \| u_x \|^2 + \beta \| \varphi - u \|^2 + \frac{\kappa}{2} \| \varphi_x + \psi \|^2 + b \| \psi_x \|^2 + \frac{m^2 c_p}{2\kappa} \| \theta_x \|^2 \\
		&\geq c \left\|
		\begin{pmatrix}u\\ \varphi \\ \psi \\ \theta\end{pmatrix}
		\right\|_V^2, 	
	\end{align*}
	where $c$ is a positive constant. Therefore, $\mathcal{L}$ is coercive. By the Lax-Milgram Theorem, system \eqref{eq:resolvent-system-1-A-1} has a unique solution $(u, \varphi, \psi, \theta) \in V$.
	
From equation \eqref{eq:resolvent-system-A-1}$_8$, it follows that
	\begin{align}
		q=-\frac{1}{\sigma}(\theta_x+\tau f_8) \in L^2.
	\end{align}
	From \eqref{eq:resolvent-system-A-1}$_7$, we also obtain $q_x=-m(\Phi_x+\Psi)-\rho_3 f_7 \in L^2$. Thus $q \in H^1$.
	
	Using standard elliptic regularity theory for \eqref{eq:resolvent-system-1-A-1}$_{1,2,3}$, we derive
	\begin{align}\label{jlll}
		u, \varphi, \psi \in H^2.
	\end{align}
	
	Setting $\dt \varphi = \dt \psi = \dt \theta = 0$ in \eqref{LF} yields
	\begin{align*}
		\alpha(u_x,\dt u_x) = (\beta(\varphi - u) + g_1, \dt u), \quad \forall \dt u \in H_*^1.
	\end{align*}
	Since $u\in H^2$, we can integrate by parts and rearrange the terms to obtain
	\begin{equation}\label{vb}
		(\alpha u_{xx}+\beta(\varphi - u) + g_1,\dt u)=\alpha \left(u_x(L)\dt u(L)-u_x(0)\dt u(0)\right), \quad \forall \dt u \in H_*^1.
	\end{equation}
	
	For $m$ an integer with $m\gg1$, we choose a nonnegative smooth function $\xi_m(x)$ defined on $\left[0,\frac{1}{m}\right]$ such that
	\[
	\xi_m(0)=1, \ 0\le\xi_m(x)\le 1 \text{ for } x\in\left(0,\frac{1}{m}\right), \ \xi_m\left(\frac{1}{m}\right)=0.
	\]
	Define
	\[
	\tilde{u}_m(x)=
	\begin{cases}
		\xi_m(x), & 0\le x\le\frac{1}{m},\\
		0, & \frac{1}{m}\le x\le L-\frac{1}{m},\\
		-\xi_m(L-x), & L-\frac{1}{m}\le x\le L.
	\end{cases}
	\]
	A series of straightforward calculations shows that $\tilde{u}_m(x)\in H_*^1$ and $\|\tilde{u}_m\|^2 \leq \frac{2}{m}$. Then by \eqref{vb}, we have
	\begin{equation}\label{vb1}
		(\alpha u_{xx}+\beta(\varphi - u) + g_1,\tilde{u}_m)=\alpha \left(-u_x(L)-u_x(0)\right).
	\end{equation}
	By Cauchy-Schwarz's inequality, we derive
	\[
	\left| (\alpha u_{xx}+\beta(\varphi - u) + g_1,\tilde{u}_m)\right|\le\left(\frac{2}{m}\right)^{\frac{1}{2}}\|\alpha u_{xx}+\beta(\varphi - u) + g_1\|\to0
	\]
	as $m\to\infty$.
	 Letting $m\to\infty$ in \eqref{vb1} gives
	\begin{equation}\label{vb2}
		u_x(L)+u_x(0)=0.
	\end{equation}
	Let
	\[
	\tilde{u}_m(x)=
	\begin{cases}
		\xi_m(x), & 0\le x\le \frac{1}{m},\\
		0, & \frac{1}{m}\le x\le\frac{L}{2}-\frac{1}{m},\\
		-\xi_m(\frac{L}{2}-x), &
		\frac{L}{2}-\frac{1}{m}\le x\le\frac{L}{2},\\
		-\xi_m(-\frac{L}{2}+x), &
		\frac{L}{2}\le x\le\frac{L}{2}+\frac{1}{m},\\
		0, &
		\frac{L}{2}+\frac{1}{m}\le x\le L-\frac{1}{m},\\
		\xi_m(L-x), & L-\frac{1}{m}\le x\le L.
	\end{cases}
	\]
	A series of straightforward calculations shows that $\tilde{u}_m(x)\in H_*^1$ and $\|\tilde{u}_m\|^2 \le\frac{4}{m}$. Then by \eqref{vb}, we have
	\begin{equation}\label{gvb1}
		(\alpha u_{xx}+\beta(\varphi - u) + g_1,\tilde{u}_m)=\alpha \left(u_x(L)-u_x(0)\right).
	\end{equation}
	By Cauchy-Schwarz's inequality, we derive
	\[
	\left| (\alpha u_{xx}+\beta(\varphi - u) + g_1,\tilde{u}_m)\right|\le\frac{2}{m^{\frac{1}{2}}}\|\alpha u_{xx}+\beta(\varphi - u) + g_1\|\to0
	\]
	as $m\to\infty$.
    Letting $m\to\infty$ in \eqref{gvb1} gives
	\begin{equation}\label{gvb2}
		u_x(L)-u_x(0)=0.
	\end{equation}
	
	From \eqref{vb2} and \eqref{gvb2}, it follows that
	\[
	u_x(L)=u_x(0)=0.
	\]
	Substituting back into \eqref{vb}, we obtain
	\[
	(\alpha u_{xx}+\beta(\varphi - u) + g_1,\tilde{u})=0, \quad \forall \tilde{u}\in H_*^1,
	\]
	which implies
	\[
	\alpha u_{xx}+\beta(\varphi - u) + g_1=0.
	\]
	Thus, we have
	\begin{equation}\label{jl6}
		u\in H^2\text{ with } u_x\in H_0^1 \text{ satisfies } \alpha u_{xx}+\beta(\varphi - u) + g_1=0.
	\end{equation}
	
	By similar reasoning,
	\begin{equation}\label{jl7}
		\varphi\in H^2\text{ with } \varphi_x\in H_0^1\text{ satisfies } -\kappa(\varphi_x + \psi)_x + \beta(\varphi - u) + m\theta_x = g_2.
	\end{equation}
	Consequently, $U \in \mathcal{D}(\mathcal{A}_1)$ is obtained. Clearly, the above proof shows that $\|U\|_{\mathcal{H}_1}=\|\mathcal{A}_1^{-1}F\|_{\mathcal{H}_1} \leq C\|F\|_{\mathcal{H}_1}$ for some positive constant $C$. Thus, $0 \in \varrho(\mathcal{A}_1)$.
\end{proof}

\subsection{Polynomial Decay}\label{sec2.2}
In this section, we establish the polynomial stability of the \(C_0\)-semigroup generated by operator \(\mathcal{A}_1\) (see Section 2 for its definition and domain \(\mathcal{D}(\mathcal{A}_1) \subset \mathcal{H}_1\)), which corresponds to Theorem \ref{Polynomial stability}. Our analysis is conducted by distinguishing two critical scenarios based on the wave speeds associated with the bridge deck: the case of equal wave speeds and the case of distinct wave speeds.

The proof of Theorem \ref{Polynomial stability} relies on the following classical stability criterion for \(C_0\)-semigroups on Hilbert spaces, which characterizes exponential and polynomial stability in terms of the resolvent estimate of the generator
\begin{theorem}\label{Stability Theorem}\cite{MR1681343,MR2606945}
	Let \(X\) be a Hilbert space, and let \(\{e^{\mathcal{A} t}\}_{t \geq 0}\) be a contraction \(C_0\)-semigroup on \(X\) with generator \(\mathcal{A}: \mathcal{D}(\mathcal{A}) \subset X \to X\). Suppose that the imaginary axis is contained in the resolvent set of \(\mathcal{A}\), i.e., \(i\mathbb{R} \subset \rho(\mathcal{A})\). Then the following assertions hold
	\begin{enumerate}[(i)]
		\item The semigroup \(\{e^{\mathcal{A} t}\}_{t \geq 0}\) is exponentially stable, i.e., $\|e^{\mathcal{A} t}\|_{\mathcal{L}(X \to X)} \le Me^{-\omega t}$ for some constant $M,\omega>0$, if and only if
        \[
        \limsup_{|\lambda| \to \infty} \|(i\lambda I - \mathcal{A})^{-1}\|_{\mathcal{L}(X \to X)} < \infty;
        \]
		\item  The semigroup \(\{e^{\mathcal{A} t}\}_{t \geq 0}\) is polynomially stable of order \(\frac{1}{l}\) (where \(l > 0\)), i.e., $\|e^{\mathcal{A} t}\phi\|_X\le Mt^{-\frac1l}\|\phi\|_{\mathcal{D}(\mathcal{A})}$ for $\phi\in \mathcal{D}(\mathcal{A})$ and some constant $M>0$, if and only if
        \[
        \limsup_{|\lambda| \to \infty} |\lambda|^{-l} \|(i\lambda I - \mathcal{A})^{-1}\|_{\mathcal{L}(X \to X)} < \infty.
        \]
	\end{enumerate}
\end{theorem}

  For any $F = (f_1,f_2,f_3,f_4,f_5,f_6,f_7,f_8) \in \mathcal{H}_1$, consider the resolvent equation $(i\lambda I-\mathcal{A}_1)U=F$, where $U=(u,v,\varphi,\Phi,\psi,\Psi,\theta,q) \in \mathcal{D}(\mathcal{A}_1)$. From the definition of $\mathcal{A}_1$ \eqref{definition:A-1}, it follows that
  \begin{align}\label{Resolvent equations}
  	\begin{cases}
  		i\lambda u - v=f_1   \quad &\text{in } H_*^1,\\
  		i\lambda \rho v -\alpha u_{xx}-\beta(\varphi - u)+ \gamma v =\rho f_2 \quad &\text{in } L_*^2,\\
  		i\lambda\varphi - \Phi =f_3\quad &\text{in } H_*^1,\\
  		i\lambda\rho_1 \Phi - \kappa(\varphi_{x}+\psi)_x + \beta (\varphi - u) +  m\theta_{x}=\rho_1f_4  \quad &\text{in } L_*^2,\\
  		i\lambda \psi - \Psi =f_5 \quad &\text{in } H_0^1,\\
  		i\lambda \rho_2 \Psi - b \psi_{xx} + \kappa(\varphi_{x} + \psi) - m\theta =\rho_2f_6 \quad &\text{in } L^2,\\
  		i\lambda \rho_3 \theta + q_x + m(\Phi_{x} + \Psi) =\rho_3f_7 \quad &\text{in } L^2,\\
  		i\lambda \tau q + \sigma q + \theta_x = \tau f_8 \quad &\text{in } L^2.
  	\end{cases}
  \end{align}

  \begin{lemma}\label{lemma-iR}
  	The resolvent set of the operator $\mathcal{A}_1$ contains the imaginary axis, i.e., $i\mathbb{R} \subset \rho(\mathcal{A}_1)$.
  \end{lemma}
  \begin{proof}
  	By contradiction, suppose $i\mathbb{R} \not\subset \varrho(\mathcal{A}_1)$. Since $0 \in \varrho(\mathcal{A}_1)$, there must exists a real number $\lambda \neq 0$, such that $i\lambda \not\in \varrho(\mathcal{A}_1)$, i.e., $i\lambda \in \sigma(\mathcal{A}_1)$.
  	The combination of $\mathcal{A}_1$ being a closed operator and the compact embedding of $\mathcal{D}(\mathcal{A}_1)$into
  	$\mathcal{H}_1$ implies that $\sigma(\mathcal{A}_1)$ contains only eigenvalues. Thus, $i \lambda$ is an eigenvalue and therefore corresponds to a non-zero eigenfunction $U=(u,v,\varphi,\Phi,\psi,\Psi,\theta,q)$. Therefore, $F=(i\lambda I-\mathcal{A}_1)U = 0$. \\
  	Taking the $\mathcal{H}_1$-inner product of $(i\lambda I-\mathcal{A}_1)U$ with $U$, from the proof of Lemma\ref{dissipative-1}, we obtain
  	\begin{align*}
  		0={\rm Re}(F,U)_{\mathcal{H}_1}={\rm Re}((i\lambda I-\mathcal{A}_1)U,U)_{\mathcal{H}_1}=-{\rm Re}(\mathcal{A}_1U,U)_{\mathcal{H}_1}=\gamma\|v\|^2+\sigma\|q\|^2.
  	\end{align*}
  	Thus, $v=q=0$. It is clear from equations \eqref{Resolvent equations} that $u=\varphi=\Phi=\psi=\Psi=\theta=0$. Therefore, $U=0$, contradicting $U \neq 0$. Consequently, $i\mathbb{R} \subset \rho(\mathcal{A}_1)$.
  \end{proof}
  \begin{lemma}\label{lemma-v,q}
  	The following estimates hold
  	\begin{align}
  		\|v\|^2 \leq \frac{1}{\gamma}\|U\|_{\mathcal{H}_1}\|F\|_{\mathcal{H}_1}, \label{estimate-v}\\
  		\|q\|^2 \leq \frac{1}{\sigma}\|U\|_{\mathcal{H}_1}\|F\|_{\mathcal{H}_1}. \label{estimate-q}
  	\end{align}
  \end{lemma}

  \begin{proof}
  	Take the $\mathcal{H}_1$-inner product of $(i\lambda I-\mathcal{A}_1)U=F$ with $U$ and compute the real part
  	\begin{align*}
  		\|F\|_{\mathcal{H}_1}\|U\|_{\mathcal{H}_1}
  		&\geq |\operatorname{Re}(F,U)_{\mathcal{H}_1}| = |\operatorname{Re}(i\lambda U-\mathcal{A}_1U,U)_{\mathcal{H}_1}| \\
  		&= |\operatorname{Re}\{i\lambda(U,U)_{\mathcal{H}_1} - (\mathcal{A}_1U,U)_{\mathcal{H}_1}\}| = |-\operatorname{Re}(\mathcal{A}_1U,U)_{\mathcal{H}_1}|\\
  		&= \gamma \|v\|^2 + \sigma \|q\|^2.
  	\end{align*}
  	Estimates \eqref{estimate-v} and \eqref{estimate-q} can be obtained directly.
  \end{proof}
  \begin{lemma}\label{lemma-ux}
  	There exists a constant $C>0$ such that
  	\begin{align}\label{estimate-ux-1}
  		\|u_x\|^2 \leq C(\|U\|_{\mathcal{H}_1}\|F\|_{\mathcal{H}_1} + \|v\|\|\varphi\|).
  	\end{align}
  	In particular, for any $\epsilon>0$, there exists a constant $C_{\epsilon}>0$ (independent of $\lambda$) such that
  	\begin{align}\label{estimate-ux-2}
  		\|u_x\|^2 \leq \epsilon\|U\|_{\mathcal{H}_1}^2 + C_{\epsilon}\|F\|_{\mathcal{H}_1}^2
  	\end{align}
  	for $|\lambda| \geq 1$.
  \end{lemma}
  \begin{proof}
  	Taking the $L^2$-inner product of equation \eqref{Resolvent equations}$_2$ with $u$ and then applying integration by parts yield
  	\begin{align}\label{eq2.1}
  		i \lambda \rho (v, u) + \alpha \| u_x \|^2 - \beta (\varphi, u) + \beta \| u \|^2 + \gamma (v, u) = \rho (f_2, u).
  	\end{align}
  	By \eqref{Resolvent equations}$_1$, we get
  	\begin{align}
  		&i \lambda \rho (v, u) = - \rho(v,i\lambda u) = - \rho(v,v+f_1) = -\rho \|v\|^2 - \rho (v,f_1),\label{eq2.2}\\
  		&\beta(\varphi,u) = \beta(\varphi,\frac{1}{i\lambda}(v+f_1))=\frac{i\beta}{\lambda}(\varphi,v)+\frac{i\beta}{\lambda}(\varphi,f_1).\label{eq2.3}
  	\end{align}
  	Substitute \eqref{eq2.2} and \eqref{eq2.3} into \eqref{eq2.1}
  	\begin{align*}
  		\alpha \| u_x \|^2 + \beta \| u \|^2
  		= \rho \| v \|^2 + \rho (v, f_1) + \frac{i\beta}{\lambda} (\varphi, v) +\frac{i\beta}{ \lambda} (\varphi, f_1) - \gamma (v, u) + \rho (f_2, u).
  	\end{align*}
  	Thus,
  	\begin{align*}
  		\alpha \|u_x\|^2 + \beta \|u\|^2
  		\leq C\|U\|_{\mathcal{H}_1}\|F\|_{\mathcal{H}_1} + \beta \|\varphi\|\|v\| + \gamma \|v\|\|u\|,
  	\end{align*}
  	for $|\lambda| \geq 1$.
  	By Young's inequality, we get
  	\begin{align*}
  		\|u_x\|^2 \leq C(\|U\|_{\mathcal{H}_1}\|F\|_{\mathcal{H}_1} + \|v\|\|\varphi\|).
  	\end{align*}
  	By estimate \eqref{estimate-v} and Young's inequality, we obtain
  	\begin{align*}
  		\|u_x\|^2\leq C\|U\|_{\mathcal{H}_1}\|F\|_{\mathcal{H}_1} + (\epsilon \|\varphi\|^2 + C_{\epsilon}\|v\|^2) \leq \epsilon\|U\|_{\mathcal{H}_1}^2 + C_{\epsilon}\|F\|_{\mathcal{H}_1}^2.
  	\end{align*}
  \end{proof}
  \begin{lemma}\label{lemma-theta}
  	There exists a constant $C>0$ such that
  	\begin{align}\label{estimate-thete-1}
  		\|\theta\|^2 \leq C(\|U\|_{\mathcal{H}_1}\|F\|_{\mathcal{H}_1} + \|q\|\|\Psi\| + \|q\|\|\Phi\|).
  	\end{align}
  	In particular, for any $\epsilon>0$, there exists a constant $C_{\epsilon}>0$ (independent of $\lambda$) such that
  	\begin{align}\label{estimate-theta-2}
  		\|\theta\|^2 \leq \epsilon\|U\|_{\mathcal{H}_1}^2 + C_{\epsilon}\|F\|_{\mathcal{H}_1}^2,
  	\end{align}
  	for $|\lambda| \geq 1$.
  \end{lemma}
  \begin{proof}
  	For any $x \in (0,L)$, integrating equation \eqref{Resolvent equations}$_8$ over the interval $(0, x)$ yields
  	\begin{align}\label{eq2.4}
  		i \lambda \tau \int_{0}^{x} q(y)  dy + \sigma \int_{0}^{x} q(y) dy + \theta(x) = \tau \int_{0}^{x} f_8(y)  dy.
  	\end{align}
  	Take the $L^2$-inner product of equation \eqref{eq2.4} with $\theta$
  	\begin{align}\label{eq-theta}
  		R_1 + \sigma \left( \int_{0}^{x} q(y)  dy, \theta \right) + \|\theta\|^2 = \tau \left( \int_{0}^{x} f_8(y)  dy, \theta \right)
  	\end{align}
  	where $R_1$ is defined as
  	\begin{align*}
  		R_1 = i \lambda \tau \left( \int_{0}^{x} q(y)  dy, \theta \right).
  	\end{align*}
  	From equation \eqref{Resolvent equations}$_7$, we obtain
  	\begin{align}\label{eq2.5}
  		i\lambda \theta = f_7 - \frac{1}{\rho_3} q_x - \frac{m}{\rho_3}(\Phi_x + \Psi).
  	\end{align}
  	Substitute equation \eqref{eq2.5} into $R_1$
  	\begin{align*}
  		R_1 &= -\tau \left( \int_{0}^{x} q(y)  dy, i\lambda \theta \right) \\
  		&= -\tau \left( \int_{0}^{x} q(y)  dy,  f_7 - \frac{1}{\rho_3} q_x - \frac{m}{\rho_3} \Phi_x - \frac{m}{\rho_3} \Psi \right) \\
  		&= -\tau \left( \int_{0}^{x} q(y)  dy,  f_7 \right) + \frac{\tau}{\rho_3} \left( \int_{0}^{x} q(y)  dy,  q_x \right) + \frac{\tau m}{\rho_3} \left( \int_{0}^{x} q(y)  dy, \Phi_x \right) + \frac{\tau m}{\rho_3} \left( \int_{0}^{x} q(y)  dy,  \Psi \right)  \\
  		&= -\tau \left( \int_{0}^{x} q(y)  dy,  f_7 \right) + \frac{\tau m}{\rho_3} \left( \int_{0}^{x} q(y)  dy,  \Psi \right) + R_2 - \frac{\tau}{\rho_3} \|q\|^2 - \frac{\tau m}{\rho_3}(q,\Phi).
  	\end{align*}
  	where $R_2$ is given by
  	\begin{align*}
  		R_2=\frac{\tau}{\rho_3}  \int_{0}^{L} q(y)  dy  (\overline{q}(L) + m  \overline{\Phi}(L))
  	\end{align*}
  	Integrating equation \eqref{Resolvent equations}$_7$ over the interval $(x, L)$, we get
  	\begin{align}\label{eq2.6}
  		i\lambda\rho_3 \int_x^L \theta(y)  dy + q(L) - q(x) + m\Phi(L) - m\Phi(x) + m \int_x^L \Psi(y)  dy = \rho_3 \int_x^L f_7(y)  dy.
  	\end{align}
  	Multiplying equation \eqref{eq2.6} by $\int_{0}^{L}\overline{q}(x)dx$, we derive
  	\begin{align*}
  		\frac{\rho_3}{\tau}\overline{R}_2 =& \big(q(L) + m\Phi(L)\big) \int_0^L \overline{q}(x) dx \\
  		=& \big(q(x) + m\Phi(x)\big) \int_0^L \overline{q}(x) dx - R_3 - m \int_x^L \Psi(y) dy \int_0^L \overline{q}(x) dx  + \rho_3 \int_x^L f_7(y) dy \int_0^L \overline{q}(x) dx,
  	\end{align*}
  	where $R_3=i\lambda\rho_3 \int_x^L \theta(y) dy \int_0^L \overline{q}(x) dx $. Equation \eqref{Resolvent equations}$_8$ yields
  	\begin{align}\label{eq2.7}
  		i\lambda q=f_8 - \frac{\sigma}{\tau}q - \frac{1}{\tau}\theta_x.
  	\end{align}
  	From equation \eqref{eq2.7}, it follows that
  	\begin{align*}
  		R_3 &= -\rho_3 \int_x^L \theta(y)  dy \int_0^L \overline{i\lambda q(x)}  dx \\
  		&= -\rho_3 \int_x^L \theta(y)  dy \int_0^L \left( \overline{f_8} - \frac{\sigma}{\tau}\overline{q} - \frac{1}{\tau} \overline{\theta_x} \right) dx \\
  		&= -\rho_3 \int_x^L \theta(y)  dy \int_0^L \overline{f_8}(x)  dx + \frac{\sigma \rho_3}{\tau} \int_x^L \theta(y)  dy \int_0^L \overline{q}(x)  dx.
  	\end{align*}
  	Substitute $R_3$ into $R_2$, then substitute $R_2$ into $R_1$, and further substitute $R_1$ into equation \eqref{eq-theta}. Integrate the resulting equation over the interval $(0, L)$; the estimate \eqref{estimate-thete-1} can be derived via H\"older's inequality and Young's inequality immediately. By estimate \eqref{estimate-q} and Young's inequality, estimate \eqref{estimate-theta-2} is a direct consequence of estimate \eqref{estimate-thete-1}.
  \end{proof}
  \begin{lemma}\label{lemma-varphix}
  	There exists a constant $C>0$ such that
  	\begin{align}\label{estimate-varphix-1}
  		\|\varphi_x + \psi\|^2 \leq C(\|U\|_{\mathcal{H}_1}\|F\|_{\mathcal{H}_1} + \frac{1}{\lambda}\|U\|_{\mathcal{H}_1}\|F\|_{\mathcal{H}_1} + \|q\|\|\Psi\| + \|q\|\|\Phi\| + \frac{1}{\lambda}\|q\|\|\varphi - u \|).
  	\end{align}
  	In particular, for any $\epsilon>0$, there exists a constant $C_{\epsilon}>0$ (independent of $\lambda$) such that
  	\begin{align}\label{estimate-varphix-2}
  		\|\varphi_x + \psi\|^2 \leq \epsilon\|U\|_{\mathcal{H}_1}^2 + C_{\epsilon}\|F\|_{\mathcal{H}_1}^2
  	\end{align}
  	for $|\lambda| \geq 1$.
  \end{lemma}
  \begin{proof}
  	From equations \eqref{Resolvent equations}$_{3,5}$, it follows that
  	\begin{align} \label{eq2.8}
  		\Phi_{x} = i \lambda \varphi_{x} - f_{3x}, \quad \Psi = i \lambda \psi - f_{5}.
  	\end{align}
  	Substituting equations \eqref{eq2.8} into \eqref{Resolvent equations}$_7$, we obtain
  	\begin{align}\label{eq2.9}
  		i \lambda \rho_{3} \theta + q_x + i m \lambda (\varphi_{x} + \psi) = m (f_{3x} + f_{5}) + \rho_{3} f_{7}.
  	\end{align}
  	Taking the $L^2$-inner product of \eqref{eq2.9} with $\varphi_x+\psi$ and integrating by parts, we derive
  	\begin{align}\label{eq-varphix}
  		i \lambda \rho_3 (\theta, \varphi_x + \psi) - R_4 + i m \lambda \|\varphi_x + \psi\|^2 = (m(f_{3x} + f_5) + \rho_3 f_7, \varphi_x + \psi).
  	\end{align}
  	where $R_4=(q, (\varphi_x + \psi)_x)$.
  	Equation \eqref{Resolvent equations}$_4$ yields
  	\begin{align}\label{eq2.10}
  		\kappa (\varphi_x + \psi)_x = i \lambda \rho_1 \Phi + \beta (\varphi - u) + m \theta_x - \rho_1 f_4.
  	\end{align}
  	Substitute equation \eqref{eq2.10} into $R_4$
  	\begin{align*}
  		R_4 &= \frac{1}{\kappa}(q, i \lambda \rho_1 \Phi + \beta (\varphi - u) + m \theta_x - \rho_1 f_4) \\
  		&= -\frac{i \lambda \rho_1}{\kappa}(q, \Phi) + \frac{\beta}{\kappa}(q, \varphi - u) + R_5 - \frac{\rho_1}{\kappa}(q, f_4).
  	\end{align*}
  	where $R_5=\frac{m}{\kappa}(q, \theta_x)$.
  	From equation \eqref{Resolvent equations}$_8$, it follows that
  	\begin{align}\label{eq2.11}
  		\theta_x = \tau f_8 - i\lambda \tau q - \sigma q.
  	\end{align}
  	Substitute equation \eqref{eq2.11} into $R_5$
  	\begin{align*}
  		R_5 = \frac{m}{\kappa} (q, \tau f_8 - i \lambda \tau q - \sigma q) = \frac{m \tau}{\kappa} (q, f_8) + \frac{i m \lambda \tau}{\kappa} \| q \|^2 - \frac{m \sigma}{\kappa} \| q \|^2.
  	\end{align*}
  	Substitute $R_5$ into $R_4$, and substitute $R_4$ into equation \eqref{eq-varphix}. By H\"older's inequality, we derive
  	\begin{align*}
  		\|\varphi_x + \psi\|^2 &\leq C\left(\|\theta\|\|\varphi_x + \psi\| + \frac{1}{\lambda}\|U\|_{\mathcal{H}_1}\|F\|_{\mathcal{H}_1} + \|q\|\|\Phi\| + \frac{1}{\lambda}\|q\|\|\varphi - u\| + \|q\|^2 + \frac{1}{\lambda}\|q\|^2\right) \\
  		&\leq C\left(\|U\|_{\mathcal{H}_1}\|F\|_{\mathcal{H}_1} + \frac{1}{\lambda}\|U\|_{\mathcal{H}_1}\|F\|_{\mathcal{H}_1} + \|\theta\|\|\varphi_x + \psi\| + \|q\|\|\Phi\| + \frac{1}{\lambda}\|q\|\|\varphi - u\|\right).
  	\end{align*}
  	Using estimate \eqref{estimate-thete-1}, by Young's inequality, we obtain
  	\begin{align*}
  		\|\varphi_x + \psi\|^2 &\leq C\left(\|U\|_{\mathcal{H}_1}\|F\|_{\mathcal{H}_1} + \frac{1}{\lambda}\|U\|_{\mathcal{H}_1}\|F\|_{\mathcal{H}_1} + \|\theta\|^2 + \|q\|\|\Phi\| + \frac{1}{\lambda}\|q\|\|\varphi - u\|\right)\\
  		&\leq C\left(\|U\|_{\mathcal{H}_1}\|F\|_{\mathcal{H}_1} + \frac{1}{\lambda}\|U\|_{\mathcal{H}_1}\|F\|_{\mathcal{H}_1} + \|q\|\|\Psi\| + \|q\|\|\Phi\| + \frac{1}{\lambda}\|q\|\|\varphi - u\|\right).
  	\end{align*}
  	The proof of \eqref{estimate-varphix-1}is completed. Using Young's inequality and estimate \eqref{estimate-q}, estimate \eqref{estimate-varphix-2} is a direct consequence of estimate \eqref{estimate-varphix-1}, for $|\lambda| \geq 1$.
  \end{proof}

  \begin{lemma}\label{lemma-varphi-u}
  	There exists a constant $C>0$ such that
  	\begin{align}\label{estimate-varphi-u-1}
  		\|\varphi - u\|^2 \leq C(\|U\|_{\mathcal{H}_1}\|F\|_{\mathcal{H}_1}  + \|v|\|\Phi\| + \|v\|\|\varphi\| + \|u_x\|\|\varphi_x\|).
  	\end{align}
  	In particular, for any $\epsilon>0$, there exists a constant $C_{\epsilon}>0$ (independent of $\lambda$) such that
  	\begin{align}\label{estimate-varphi-u-2}
  		\|\varphi - u\|^2 \leq \epsilon\|U\|_{\mathcal{H}_1}^2 + C_{\epsilon}\|F\|_{\mathcal{H}_1}^2
  	\end{align}
  	for $|\lambda| \geq 1$.
  \end{lemma}
  \begin{proof}
  	Take the $L^2$-inner product of equation \eqref{Resolvent equations}$_2$ with $\varphi-u$ and integrate by parts
  	\begin{align}\label{eq-varphi-u}
  		R_6 + \alpha (u_x, \varphi_x - u_x) - \beta \|\varphi - u\|^2 + \gamma (v, \varphi - u) = \rho (f_2, \varphi - u).
  	\end{align}
  	where $R_6$ is defined as
  	\begin{align*}
  		R_6=i \lambda \rho (v, \varphi - u).
  	\end{align*}
  	Using equations \eqref{Resolvent equations}$_{1,3}$, we obtain
  	\begin{align*}
  		R_6 &= -\rho (v, i \lambda \varphi - i \lambda u)
  		= -\rho (v, \Phi + f_3 - v - f_1)
  		= -\rho (v, \Phi) + \rho \|v\|^2 - \rho (v, f_3 - f_1).
  	\end{align*}
  	Substitute $R_6$ into equation \eqref{eq-varphi-u}. By H\"older's inequality, we derive
  	\begin{align*}
  		\|\varphi - u\|^2 \leq C(\|v\|\|\Phi\| + \|U\|_{\mathcal{H}_1}\|F\|_{\mathcal{H}_1} + \|u_x\|\|\varphi_x\| + \|u_x\|^2 + \|v\|\|\varphi - u\|)
  	\end{align*}
  	Using estimates \eqref{estimate-v} and \eqref{estimate-ux-1}, by Young's inequality, we obtain \eqref{estimate-varphi-u-1} immediately. Using estimates \eqref{estimate-v} and \eqref{estimate-ux-2}, \eqref{estimate-varphi-u-2} is a direct consequence of \eqref{estimate-varphi-u-1}.
  \end{proof}

  \begin{lemma}\label{lemma-Phi}
  	There exists a constant $C>0$ such that
  	\begin{align}\label{estimate-Phi-1}
  		\|\Phi\|^2 \leq C(\|U\|_{\mathcal{H}_1}\|F\|_{\mathcal{H}_1}  + \|\varphi_x+\psi|\|\varphi_x\| + \|\varphi-u\|\|\varphi\| + \|\theta\|\|\varphi_x\|).
  	\end{align}
  	In particular, for any $\epsilon>0$, there exists a constant $C_{\epsilon}>0$ (independent of $\lambda$) such that
  	\begin{align}\label{estimate-Phi-2}
  		\|\Phi\|^2 \leq \epsilon\|U\|_{\mathcal{H}_1}^2 + C_{\epsilon}\|F\|_{\mathcal{H}_1}^2
  	\end{align}
  	for $|\lambda| \geq 1$.
  \end{lemma}
  \begin{proof}
  	Take the $L^2$-inner product of equation \eqref{Resolvent equations}$_4$ with $\varphi$ and integrate by parts
  	\begin{align}\label{eq-Phi}
  		R_7 + \kappa (\varphi_x + \psi, \varphi_x) + \beta (\varphi - u, \varphi) - m (\theta, \varphi_x) = \rho_1 (f_4, \varphi),
  	\end{align}
  	where $R_7 = i \lambda \rho_1 (\Phi, \varphi)$. From equation \eqref{Resolvent equations}$_3$, it follows that
  	\begin{align}
  		R_7 = -\rho_1 (\Phi, i \lambda \varphi) = -\rho_1 (\Phi, \Phi + f_3) = -\rho_1 \|\Phi\|^2 - \rho_1 (\Phi, f_3).
  	\end{align}
  	Substitute $R_7$ into equation \eqref{eq-Phi}. Then by H\"older's inequality, we get \eqref{estimate-Phi-1} immediately. Using estimates \eqref{estimate-varphix-2}, \eqref{estimate-varphi-u-2} and \eqref{estimate-theta-2}, \eqref{estimate-Phi-2} is a direct consequence of \eqref{estimate-Phi-1}.
  \end{proof}

  \begin{lemma}\label{lemma-psix}
  	For $|\lambda| \geq \max\{1, \sigma\}$, the following inequalities hold
  	\begin{enumerate}[(i)]
  		\item if $\frac{\rho_1}{\kappa}=\frac{\rho_2}{b}$, for any $\epsilon>0$, there exists a constant $C_{\epsilon}>0$ (independent of $\lambda$) such that
  		\begin{align}
  			\|\psi_x\|^2 \leq \epsilon\|U\|_{\mathcal{H}_1}^2 + C_{\epsilon} \lambda ^4 \|F\|_{\mathcal{H}_1}^2;
  		\end{align}
  		\item if $\frac{\rho_1}{\kappa} \neq \frac{\rho_2}{b}$, for any $\epsilon>0$, there exists a constant $C_{\epsilon}>0$ (independent of $\lambda$) such that
  		\begin{align}
  			\|\psi_x\|^2 \leq \epsilon\|U\|_{\mathcal{H}_1}^2 + C_{\epsilon}\lambda ^8\|F\|_{\mathcal{H}_1}^2.
  		\end{align}
  	\end{enumerate}
  \end{lemma}
  \begin{proof}
  	From equations \eqref{Resolvent equations}$_{3,5}$, it follows that
  	\begin{align}\label{eq2.12}
  		i \lambda \varphi_x - \Phi_x = f_{3x}, \quad i \lambda \psi - \Psi = f_5.
  	\end{align}
  	Add the two equations in \eqref{eq2.12}
  	\begin{align}\label{eq2.13}
  		i \lambda (\varphi_x + \psi) - (\Phi_x + \Psi) = f_{3 x} + f_5.
  	\end{align}
  	Taking the $L^2$-inner product of equation \eqref{eq2.13} with $\Psi$ and rearranging the terms, we derive
  	\begin{align}\label{eq2.14}
  		-i \lambda (\varphi_x + \psi, \Psi) - R_8 + R_9 = -(f_{3 x}, \Psi) - (f_5, \Psi).
  	\end{align}
  	where $R_8$ and $R_9$ are defined as:	$R_8 = (\Phi, \Psi_x), \quad R_9 = (\Psi, \Psi)$.
  	By equation \eqref{Resolvent equations}$_5$, substitute $\Psi_x$ and $\Psi$ into $R_8$ and $R_9$ to obtain
  	\begin{align*}
  		R_8 &= (\Phi, i \lambda \psi_x - f_{5 x}) = -i \lambda (\Phi, \psi_x) - (\Phi, f_{5 x}), \\
  		R_9 &= (\Psi, i \lambda \psi - f_5) = -i \lambda (\Psi, \psi) - (\Psi, f_5).
  	\end{align*}
  	Substitute $R_8$ and $R_9$ into equation \eqref{eq2.14}
  	\begin{align}\label{eq2.15}
  		-i \lambda (\varphi_x + \psi, \Psi) + i \lambda (\Phi, \psi_x)   -i \lambda (\Psi, \psi)  = -(f_{3 x}, \Psi) - (f_5, \Psi) - (\Phi, f_{5 x}) + (\Psi, f_5).
  	\end{align}
  	Taking the $L^2$-inner product of equation \eqref{Resolvent equations}$_4$ with $\frac{1}{\rho_1}\psi_x$ and the $L^2$-inner product of equation \eqref{Resolvent equations}$_6$ with $\frac{1}{\rho_2}\psi$, rearranging the terms, we derive
  	\begin{align}
  		i\lambda(\Phi,\psi_{x}) &= -\frac{\kappa}{\rho_{1}}(\varphi_{x} + \psi, \psi_{xx}) - \frac{\beta}{\rho_{1}}(\varphi - u, \psi_{x}) - \frac{m}{\rho_{1}}(\theta_{x}, \psi_{x}) + (f_{4}, \psi_{x}), \label{a1}\\
  		i\lambda(\Psi,\psi) &= -\frac{b}{\rho_{2}}\|\psi_{x}\|^{2} - \frac{\kappa}{\rho_{2}}(\varphi_{x} + \psi, \psi) + \frac{m}{\rho_{2}}(\theta, \psi) + (f_{6}, \psi). \label{a2}
  	\end{align}
  	Substitute the expressions for $i\lambda(\Phi,\psi_{x})$ and $i\lambda(\Psi,\psi)$ \eqref{a1} and
  	\eqref{a2} into equation \eqref{eq2.15}
  	\begin{align}\label{l1}
  		&-i \lambda (\varphi_x + \psi, \Psi) -\frac{\kappa}{\rho_{1}}(\varphi_{x} + \psi, \psi_{xx}) - \frac{\beta}{\rho_{1}}(\varphi - u, \psi_{x}) - \frac{m}{\rho_{1}}(\theta_{x}, \psi_{x}) +\frac{b}{\rho_{2}}\|\psi_{x}\|^{2} + \frac{\kappa}{\rho_{2}}(\varphi_{x} + \psi, \psi) - \frac{m}{\rho_{2}}(\theta, \psi)\notag\\
  		&= -(f_{3 x}, \Psi) - (f_5, \Psi) - (\Phi, f_{5 x}) + (\Psi, f_5) - (f_{4}, \psi_{x}) + (f_{6}, \psi).
  	\end{align}
  	Equation \eqref{Resolvent equations}$_6$ gives the following equation
  	\begin{align*}
  		-\frac{\kappa}{b \rho_1} ( \varphi_x + \psi, i \lambda \rho_2 \Psi - b \psi_{xx} + \kappa (\varphi_x + \psi) - m \theta ) = -\frac{\kappa}{b \rho_1} ( \varphi_x + \psi, \rho_2 f_6).
  	\end{align*}
  	Namely,
  	\begin{align}\label{l2}
  		i \lambda \frac{\kappa \rho_2}{b \rho_1} ( \varphi_x + \psi, \Psi ) + \frac{\kappa}{\rho_1} ( \varphi_x + \psi, \psi_{xx}) - \frac{\kappa ^2}{b \rho_1} \| \varphi_x + \psi \|^2
  		+ \frac{\kappa m}{b \rho_1} ( \varphi_x + \psi, \theta ) = -\frac{\kappa \rho_2}{b \rho_1} ( \varphi_x + \psi, f_6).
  	\end{align}
  	Summing equation \eqref{l1} and equation \eqref{l2}, we have
  	\begin{align}\label{l3}
  		&i \lambda (\frac{\kappa \rho_2}{b \rho_1} - 1) (\varphi_x + \psi, \Psi)  - \frac{\beta}{\rho_{1}}(\varphi - u, \psi_{x}) - \frac{m}{\rho_{1}}(\theta_{x}, \psi_{x}) +\frac{b}{\rho_{2}}\|\psi_{x}\|^{2} + \frac{\kappa}{\rho_{2}}(\varphi_{x} + \psi, \psi) - \frac{m}{\rho_{2}}(\theta, \psi) - \frac{\kappa ^2}{b \rho_1} \| \varphi_x + \psi \|^2 \notag\\
  		& + \frac{\kappa m}{b \rho_1} ( \varphi_x + \psi, \theta )
  		= -(f_{3 x}, \Psi) - (f_5, \Psi) - (\Phi, f_{5 x}) + (\Psi, f_5) - (f_{4}, \psi_{x}) + (f_{6}, \psi) -\frac{\kappa \rho_2}{b \rho_1} ( \varphi_x + \psi, f_6).
  	\end{align}
  	\paragraph{Case 1:}If $\frac{\rho_1}{\kappa}=\frac{\rho_2}{b}$, i.e., $\frac{\kappa \rho_2}{b \rho_1} = 1$.
  	From equation \eqref{l3}, it follows that
  	\begin{align*}
  		\|\psi_x\|^2 \leq C \left(\|\varphi-u\|^2 + \|\theta_x\|^2 + \|\varphi_x+\psi\|^2 + \|U\|_{\mathcal{H}_1} \|F\|_{\mathcal{H}_1}  \right).
  	\end{align*}
  	Now, we estimate $\|\theta_x\|^2$. By equation \eqref{Resolvent equations}$_8$, we derive
  	\begin{align}\label{estimate-thetax}
  		\|\theta_x\|^2 &= \|\tau f_8 - (i \lambda \tau + \sigma) q\|^2 \leq 2 \tau^2 \|f_8\|^2 + C \lambda ^2 \|q\|^2 \leq 2 \tau^2 \|F\|_{\mathcal{H}_1}^2 + C \lambda^2 \|U\|_{\mathcal{H}_1} \|F\|_{\mathcal{H}_1} \notag\\
  		&\leq 2 \tau^2 \|F\|_{\mathcal{H}_1}^2 + C \lambda ^2 \left( \frac{\epsilon}{\lambda ^2} \|U\|_{\mathcal{H}_1}^2 + C_{\epsilon} \lambda^2 \|F\|_{\mathcal{H}_1}^2 \right) \leq \epsilon \|U\|_{\mathcal{H}_1}^2 + C_{\epsilon} \lambda^4 \|F\|_{\mathcal{H}_1}^2,
  	\end{align}
  	for $|\lambda| \geq \sigma$.
  	From estimates \eqref{estimate-thetax}, \eqref{estimate-varphi-u-2} and \eqref{estimate-varphix-2}, it follows that
  	\begin{align}
  		\|\psi_x\|^2 \leq \epsilon \|U\|_{\mathcal{H}_1}^2 + C_{\epsilon} \lambda^4 \|F\|_{\mathcal{H}_1}^2,
  	\end{align}
  	for $|\lambda| \geq \max\{1, \sigma\}$.
  	\paragraph{Case 2:}If $\frac{\rho_1}{\kappa}\neq\frac{\rho_2}{b}$, i.e., $\frac{\kappa \rho_2}{b \rho_1} \neq 1$.
  	From equation \eqref{l3}, it follows that
  	\begin{align*}
  		\|\psi_x\|^2 \leq C \left(\lambda\|\varphi_x + \psi\|\|\Psi\| + \|\varphi-u\|^2 + \|\theta_x\|^2 + \|\varphi_x+\psi\|^2 + \|U\|_{\mathcal{H}_1} \|F\|_{\mathcal{H}_1}  \right).
  	\end{align*}
  	By Young's inequality, estimate \eqref{estimate-varphix-1} yields
  	\begin{align*}
  		\|\varphi_x + \psi\|^2 \leq C\left(\epsilon \|U\|_{\mathcal{H}_1}^2 + \frac{1}{\epsilon}\|U\|_{\mathcal{H}_1}\|F\|_{\mathcal{H}_1}\right).
  	\end{align*}
  	Consequently,
  	\begin{align}\label{estimate1}
  		\lambda\|\varphi_x + \psi\|\|\Psi\| &\leq \frac{\epsilon}{2}\|\Psi\|^2 + \frac{\lambda ^2}{2\epsilon} \|\varphi_x + \psi\|^2
  		\leq \epsilon \|U\|_{\mathcal{H}_1}^2 + \frac{C \lambda ^2}{2\epsilon}\left(\frac{\epsilon^2}{\lambda^2}\|U\|_{\mathcal{H}_1}^2 + \frac{\lambda^2}{\epsilon^2}\|U\|_{\mathcal{H}_1}\|F\|_{\mathcal{H}_1}\right) \notag\\
  		&\leq \epsilon \|U\|_{\mathcal{H}_1}^2 + \frac{C \lambda ^4}{\epsilon^3}\|U\|_{\mathcal{H}_1}\|F\|_{\mathcal{H}_1}
  		\leq \epsilon \|U\|_{\mathcal{H}_1}^2 + \frac{C \lambda ^4}{\epsilon^3}\left(\frac{\epsilon^4}{\lambda^4}\|U\|_{\mathcal{H}_1}^2 + \frac{\lambda^4}{\epsilon^4}\|F\|_{\mathcal{H}_1}^2\right) \notag\\
  		&\leq \epsilon \|U\|_{\mathcal{H}_1}^2 + \frac{C \lambda^8}{\epsilon^7}\|F\|_{\mathcal{H}_1}^2
  		\leq \epsilon \|U\|_{\mathcal{H}_1}^2 + C_{\epsilon}\lambda^8 \|F\|_{\mathcal{H}_1}^2.
  	\end{align}
  	From estimates \eqref{estimate1}, \eqref{estimate-thetax}, \eqref{estimate-varphi-u-2} and \eqref{estimate-varphix-2}, it follows that
  	\begin{align*}
  		\|\psi_x\|^2 \leq \epsilon \|U\|_{\mathcal{H}_1}^2 + C_{\epsilon} \lambda^8 \|F\|_{\mathcal{H}_1}^2,
  	\end{align*}
  	for $|\lambda| \geq \max\{1, \sigma\}$.
  \end{proof}

  \begin{lemma}\label{lemma-Psi}
  	For $|\lambda| \geq \max\{1, \sigma\}$, the following inequalities hold
  	\begin{enumerate}[(i)]
  		\item if $\frac{\rho_1}{\kappa}=\frac{\rho_2}{b}$, for any $\epsilon>0$, there exists a constant $C_{\epsilon}>0$ (independent of $\lambda$) such that
  		\begin{align}
  			\|\Psi\|^2 \leq \epsilon\|U\|_{\mathcal{H}_1}^2 + C_{\epsilon} \lambda ^4 \|F\|_{\mathcal{H}_1}^2;
  		\end{align}
  		\item if $\frac{\rho_1}{\kappa} \neq \frac{\rho_2}{b}$, for any $\epsilon>0$, there exists a constant $C_{\epsilon}>0$ (independent of $\lambda$) such that
  		\begin{align}
  			\|\Psi\|^2 \leq \epsilon\|U\|_{\mathcal{H}_1}^2 + C_{\epsilon}\lambda ^8\|F\|_{\mathcal{H}_1}^2.
  		\end{align}
  	\end{enumerate}
  \end{lemma}
  \begin{proof}
  	Take the $L^2$-inner product of equation \eqref{Resolvent equations}$_6$ with $\psi$, and use equation \eqref{Resolvent equations}$_5$
  	\begin{align*}
  		-\rho_2\|\Psi\|^2 - \rho_2(\Psi,f_5) + b\|\psi_x\|^2 + \kappa(\varphi_x+\psi,\psi) - m(\theta,\psi) = \rho_2(f_6,\psi).
  	\end{align*}
  	Therefore, by H\"older's inequality,
  	\begin{align*}
  		\|\Psi\|^2 \leq C(\|\psi_x\|^2+\|\varphi_x+\psi\|^2+\|\theta\|^2+\|U\|_{\mathcal{H}_1}\|F\|_{\mathcal{H}_1}).
  	\end{align*}
  	Consequently, Lemma \ref{lemma-Psi} is a direct consequence of Lemmas \ref{lemma-theta}, \ref{lemma-varphix} and \ref{lemma-psix}.
  \end{proof}
  \begin{proof}[Proof of Theorem \ref{Polynomial stability}]
  From Lemmas \ref{lemma-v,q}-\ref{lemma-Psi}, for $|\lambda| \geq \max\{1, \sigma\}$, the following estimates hold
  \paragraph{Case 1:} if $\frac{\rho_1}{\kappa}=\frac{\rho_2}{b}$, there exists a constant $C> 0$ (independent of $\lambda$) such that $\|U\|_{\mathcal{H}_1} \leq C \lambda^2\|F\|_{\mathcal{H}_1}$;
  \paragraph{Case 2:} if $\frac{\rho_1}{\kappa} \neq \frac{\rho_2}{b}$, there exists a constant $C> 0$ (independent of $\lambda$) such that $\|U\|_{\mathcal{H}_1} \leq C \lambda^4\|F\|_{\mathcal{H}_1}$.\\
  Combined with Lemma \ref{lemma-iR}, by Theorem
  \ref{Stability Theorem}, we conclude that Theorem \ref{Polynomial stability} holds.
  \end{proof}

  \subsection{Non-exponential Stability}\label{sec2.3}
  In this section, we demonstrate that the system fails to be exponentially stable under certain coefficient conditions and give the proof of Theorem \ref{Non-exponential Stability}.
  \begin{proof}[Proof of Theorem \ref{Non-exponential Stability}]
  We employ the strategy of demonstrating that the operators $(i\lambda I-\mathcal{A}_1)^{-1}$ are not uniformly bounded. Specifically, we will show that there exists a sequence of real numbers $\lambda_n$ with $|\lambda_n| \to \infty$, and a sequence of functions $F_n \in \mathcal{H}_1$ such that \{$F_n$\} is bounded in $\mathcal{H}_1$, satisfying the equation $(i\lambda_n I-\mathcal{A}_1)U_n = F_n$, where $U_n \in \mathcal{D}(\mathcal{A}_1)$ and $\|U_n\|_{\mathcal{H}_1} \to \infty$, $n \in \mathbb{N}$.
  	
  	We define $F_n=(0,0,0,0,0,\rho_2 ^{-1}sin\omega \lambda_n x,0,0) \in \mathcal{H}_1$, $\lambda_n = \frac{n\pi}{\omega L}$, where $\omega=\sqrt{\frac{\rho_2}{b}}$, $n \in \mathbb{N}$. For notational simplicity, we will omit the subscript $n$ in what follows. From the definition of $\mathcal{A}_1$, $(i\lambda_n I-\mathcal{A}_1)U_n = F_n$ is equivalent to the following equations
  	\begin{align}\label{j1}
  		\begin{cases}
  			i\lambda u - v=0,\\
  			i\lambda \rho v -\alpha u_{xx}-\beta(\varphi - u)+ \gamma v =0,\\
  			i\lambda\varphi - \Phi =0,\\
  			i\lambda\rho_1 \Phi - \kappa(\varphi_{x}+\psi)_x + \beta (\varphi - u) +  m\theta_{x}=0  ,\\
  			i\lambda \psi - \Psi =0 ,\\
  			i\lambda \rho_2 \Psi - b \psi_{xx} + \kappa(\varphi_{x} + \psi) - m\theta =sin\omega\lambda x ,\\
  			i\lambda \rho_3 \theta + q_x + m(\Phi_{x} + \Psi) =0 ,\\
  			i\lambda \tau q + \sigma q + \theta_x = 0 .
  		\end{cases}
  	\end{align}
  	Equations \eqref{j1}$_{1,3,5}$ yield
  	\begin{align}\label{j2}
  		 v = i \lambda u, \quad \Phi = i \lambda \varphi, \quad \Psi = i \lambda \psi .
  	\end{align}
  	Substitute equations \eqref{j2} into equations \eqref{j1}$_{2,4,6,7,8}$
  	\begin{align}\label{j3}
  		\begin{cases}
  			& (i \lambda \gamma - \lambda^2 \rho + \beta) u - \alpha u_{xx} - \beta \varphi = 0, \\
  			&(\beta - \lambda^2 \rho_1) \varphi - \kappa \varphi_{xx} - \kappa \psi_x - \beta u + m \theta_x = 0, \\
  			&(\kappa - \lambda^2 \rho_2) \psi - b \psi_{xx} + \kappa \varphi_x - m \theta = \sin \omega \lambda x,\\
  			& i \lambda \rho_3 \theta + q_x + i\lambda m \varphi_x + i\lambda m\psi = 0, \\
  			& i\lambda \tau q + \sigma q + \theta_x = 0.
  		\end{cases}
  	\end{align}
  	Given that $U \in \mathcal{D}(\mathcal{A}_1)$, we assume the solution of the equations \eqref{j3} takes the following form
  	\begin{align}\label{j4}
  		u=Acos\omega\lambda x, \varphi=Bcos\omega\lambda x, \psi=Csin\omega\lambda x, \theta=Dsin\omega\lambda x, q=Ecos\omega\lambda x.
  	\end{align}
  	Substituting \eqref{j4} into equations \eqref{j3}, we derive
  	\begin{align}\label{j5}
  		\begin{cases}
  			(i \lambda \gamma - \rho \lambda^2 + \beta + \alpha \omega^2 \lambda^2) A - \beta B = 0, \\
  			-\beta A + (\beta - \rho_1 \lambda^2 + \kappa \omega^2 \lambda^2) B - \kappa \omega \lambda C + m \omega \lambda D = 0, \\
  			-\kappa \omega \lambda B + (\kappa - \rho_2 \lambda^2 + b \omega^2 \lambda^2) C - m D = 1, \\
  			-im \omega \lambda B + im C + i \rho_3 D - \omega E = 0, \\
  			\omega \lambda D + (i \lambda \tau + \sigma) E = 0.
  		\end{cases}
  	\end{align}
  	From equations \eqref{j5}$_{1,5}$, we obtain
  	\begin{align}\label{j6}
  		A=\frac{\beta}{i \lambda \gamma - \rho \lambda^2 + \beta + \alpha \omega^2 \lambda^2}B, \quad E=-\frac{\omega \lambda}{i \lambda \tau + \sigma}D.
  	\end{align}
  	Substituting \eqref{j6} into equations \eqref{j5}$_{2,3,4}$, it follows that
  	\begin{align}
  		\underbrace{
  			\left[ \begin{array}{ccc}
  				q_{1}(\lambda) & -\kappa\omega\lambda & m\omega\lambda \\
  				-\kappa\omega\lambda & q_{2}(\lambda) & -m \\
  				-im\omega\lambda & im & q_{3}(\lambda)
  			\end{array} \right]
  		}_{:=M}
  		\left[ \begin{array}{c}
  			B \\
  			C \\
  			D
  		\end{array} \right]
  		=
  		\left[ \begin{array}{c}
  			0 \\
  			1 \\
  			0
  		\end{array} \right],
  	\end{align}
  	where $q_1(\lambda)$, $q_2(\lambda)$, $q_3(\lambda)$ are defined as
  	\begin{align}
  		\begin{cases}
  			q_1(\lambda) = -\dfrac{\beta^2}{i \lambda \gamma - \rho \lambda^2 + \beta + \alpha \omega^2 \lambda^2} + \beta - \rho_1 \lambda^2 + \kappa \omega^2 \lambda^2 = -\dfrac{\beta^2}{i \lambda \gamma - \rho \lambda^2 + \beta + \alpha \omega^2 \lambda^2} + \beta + \kappa\left(\dfrac{\rho_2}{b}-\dfrac{\rho_1}{\kappa}\right) \lambda^2,\\
  			q_2(\lambda) = \kappa - \rho_2 \lambda^2 + b \omega^2 \lambda^2 = \kappa,\\
  			q_3(\lambda) = i \rho_3 + \dfrac{\omega^2 \lambda}{i \lambda \tau + \sigma} = i\left(\rho_3 - \dfrac{\tau \omega ^2 \lambda^2}{\sigma^2 + \tau^2 \lambda^2}\right) + \dfrac{\sigma \omega^2 \lambda}{\sigma^2 + \tau^2 \lambda^2}.
  		\end{cases}
  	\end{align}
  	Therefore, the determinant of matrix $M$ is
  	\begin{align*}
  		\det M &= -im\omega\lambda \left( \kappa\omega m\lambda - \kappa\omega m\lambda \right) - im(-mq_1(\lambda)+\kappa m\omega^2\lambda^2)  + q_3(\lambda)(q_1(\lambda) q_2(\lambda) - \kappa^2 \omega^2\lambda^2) \\
  		&= - im(-mq_1(\lambda)+\kappa m\omega^2\lambda^2)  + q_3(\lambda)(q_1(\lambda) q_2(\lambda) - \kappa^2 \omega^2\lambda^2) \\
  		&= -im\left( \frac{m\beta^2}{i\lambda \gamma-\rho\lambda^2+\beta+\alpha\omega^2\lambda^2}  - m\beta + m \rho_1 \lambda^2\right) \\
  		&\quad + \left(i\left(\rho_3 - \dfrac{\tau \omega ^2 \lambda^2}{\sigma^2 + \tau^2 \lambda^2}\right) + \dfrac{\sigma \omega^2 \lambda}{\sigma^2 + \tau^2 \lambda^2}\right)\left( -\frac{\kappa\beta^2}{i\lambda \gamma-\rho\lambda^2+\beta+\alpha \omega^2\lambda^2}+\kappa\beta-\kappa\rho_1 \lambda^2\right).
  	\end{align*}
  	Considering the real part of $\det M$, it can be concluded that $\det M \neq 0$.
  	Thus, $C$ can be solved as
  	\begin{align}\label{C}
  		C = \dfrac{q_1(\lambda)q_3(\lambda) + im^2 \omega^2 \lambda^2}{\det M}.
  	\end{align}
  	\paragraph{Case 1:}If $\chi_0 \neq 0$, and $\chi_1 \neq 0$, as $n \to \infty$, we obtain
  	\begin{align*}
  		q_1(\lambda)q_3(\lambda) + im^2\omega^2\lambda^2 &\sim \kappa\left(\frac{\rho_2}{b}-\frac{\rho_1}{\kappa}\right)\lambda^2 \cdot i\left(\rho_3-\frac{\omega^2}{\tau}\right) + im^2\omega^2\lambda^2 \\
  		&= i\left[\kappa\left(\frac{\rho_2}{b}-\frac{\rho_1}{\kappa}\right)\left(\rho_3-\frac{\omega^2}{\tau}\right) + m^2\omega^2\right]\lambda^2 \\
  		&= i\chi_0\lambda^2,
  	\end{align*}
  	where $\chi_0 = \kappa\left(\dfrac{\rho_2}{b} -  \dfrac{\rho_1}{\kappa}\right)\left(\rho_3 - \dfrac{\rho_2}{b\tau}\right) + \dfrac{m^2 \rho_2}{b}$.
  	\begin{align*}
  		\det M &\sim -im^2 \rho_1 \lambda^2 + i \left(\rho_3 - \frac{\omega^2}{\tau}\right)(-\kappa \rho_1 \lambda^2) \\
  		&= i \rho_1 \left(-m^2 - \kappa \rho_3 + \frac{\kappa \omega^2}{\tau}\right) \lambda^2\\
  		&=i \rho_1 \chi_1 \lambda^2.
  	\end{align*}
  	where $\chi_1 = -m^2 - \kappa\rho_3 + \dfrac{\kappa\rho_2}{\tau b}$.
  	From \eqref{C}, it follows that
  	\begin{align*}
  		|C| \approx \left|\dfrac{\chi_0}{\rho_1 \chi_1}\right| > 0.
  	\end{align*}
  	Since $\Psi = i \lambda \psi = i \lambda C sin\omega \lambda x$, we derive
  	\begin{align*}
  		\|\Psi\|^2 = \lambda^2 C^2 \int_{0}^{L} sin^2\omega\lambda x dx = \lambda^2 C^2 \int_{0}^{L} sin^2\frac{n\pi}{L}x dx = \dfrac{LC^2\lambda^2}{2} \to \infty,
  	\end{align*}
  	which implies $\|U_n\|_{\mathcal{H}_1} \to \infty$.
  	\paragraph{Case 2:}If $\chi_0 \neq 0$, and $\chi_1 = 0$, as $n \to \infty$, we obtain
  	\begin{align*}
  		\det M \sim \dfrac{\sigma \omega^2}{\tau^2} \dfrac{1}{\lambda}(-\kappa \rho_1 \lambda^2) = -\dfrac{\kappa \rho_1 \sigma \omega^2}{\tau^2} \lambda.
  	\end{align*}
  	From \eqref{C}, it follows that
  	\begin{align*}
  		|C| \approx \left|\dfrac{\tau^2 \chi_0}{\kappa \rho_1 \sigma \omega^2}\lambda\right| \to \infty,
  	\end{align*}
  	which implies that
  	\begin{align*}
  		\|\Psi\|^2 = \lambda^2 C^2 \int_{0}^{L} sin^2\omega\lambda x dx = \lambda^2 C^2 \int_{0}^{L} sin^2\frac{n\pi}{L}x dx = \dfrac{LC^2\lambda^2}{2} \to \infty,
  	\end{align*}
  	which implies $\|U_n\|_{\mathcal{H}_1} \to \infty$.
  	\paragraph{Case 3:}If $\chi_0 = 0$, and $\chi_1 = 0$, (since $\chi_0 = \kappa\left(\dfrac{\rho_2}{b} -  \dfrac{\rho_1}{\kappa}\right)\left(\rho_3 - \dfrac{\rho_2}{b\tau}\right) + \dfrac{m^2 \rho_2}{b} = 0$, we get $\dfrac{\rho_2}{b} - \dfrac{\rho_1}{\kappa} \neq 0$),\\as $n \to \infty$, we obtain
  	\begin{align*}
	  		q_1(\lambda)q_3(\lambda) + im^2\omega^2\lambda^2 &\sim \kappa(\dfrac{\rho_2}{b} - \dfrac{\rho_1}{\kappa}) \lambda^2 \dfrac{\sigma\omega^2}{\tau^2} \dfrac{1}{\lambda} = \dfrac{\kappa\sigma\omega^2(\frac{\rho_2}{b} - \frac{\rho_1}{\kappa})}{\tau^2} \lambda \neq 0.
  	\end{align*}
  	From \eqref{C}, it follows that
  	\begin{align*}
  		|C| \approx \dfrac{1}{\rho_1}\left|\dfrac{\rho_2}{b} - \dfrac{\rho_1}{\kappa}\right| > 0,
  	\end{align*}
  	which implies $\|\Psi\|^2 \to \infty$. Therefore, $\|U_n\|_{\mathcal{H}_1} \to \infty$.
  	
  	The proof of Theorem \ref{Non-exponential Stability} is completed.
  \end{proof}
\section{Model Formulation with the Boltzmann Memory Term ($\delta = 1$)} \label{sec3}
When $\delta = 1$, the suspension bridge model \eqref{model1.1} incorporates the Boltzmann memory term, taking the following form

\begin{align}\label{eq:with memory}
	\begin{cases}
		\rho u_{tt} - \alpha u_{xx} - \beta(\varphi - u) + \gamma u_t = 0 & \text{in } (0,L) \times \mathbb{R}^+, \\
		\rho_1 \varphi_{tt} - \kappa(\varphi_x + \psi)_x + \beta (\varphi - u) + m \theta_x = 0 & \text{in } (0,L) \times \mathbb{R}^+, \\
		\rho_2 \psi_{tt} - b \psi_{xx} + \kappa(\varphi_x + \psi) + \int_0^\infty g(s) \psi_{xx}(t-s)  ds - m \theta = 0 & \text{in } (0,L) \times \mathbb{R}^+, \\
		\rho_3 \theta_t + q_x + m (\varphi_x + \psi)_t = 0 & \text{in } (0,L) \times \mathbb{R}^+, \\
		\tau q_t + \sigma q + \theta_x = 0& \text{in } (0,L) \times \mathbb{R}^+.
	\end{cases}
\end{align}

For the memory kernel $g(s)$ in the system, we impose the following standard assumptions, which are conventional in viscoelasticity research

\begin{assumption}\label{assumg}
 The memory kernel $g \in L^1(\mathbb{R}^+) \cap C^1(\mathbb{R}^+)$ is a nonnegative function satisfying
 \begin{equation}\label{g}
		\dt b := b - b_0 > 0 \quad \text{with } b_0 = \int_{0}^{\infty} g(s) ds, \quad \text{and } g'(s) \leq -\kappa_1 g(s) \quad \text{for all } s\in\mathbb{R}^+
	\end{equation}
	where $\kappa_1 > 0$ is a given decay rate.
\end{assumption}

Define the relative history function
\begin{equation*}
	\eta^t(x,s) := \psi(x,t) - \psi(x,t-s) \quad \text{for } t,s \geq 0.
\end{equation*}
From the relevant history function, the original system \eqref{eq:with memory} of equations is transformed into the following system
\begin{align}
	\begin{cases}
		\rho u_{tt} - \alpha u_{xx} - \beta(\varphi - u) + \gamma u_t = 0 & \text{in } (0,L) \times \mathbb{R}^+, \\
		\rho_1 \varphi_{tt} - \kappa(\varphi_x + \psi)_x + \beta (\varphi - u) + m \theta_x = 0 & \text{in } (0,L) \times \mathbb{R}^+, \\
		\rho_2 \psi_{tt} - \dt b \psi_{xx} + \kappa(\varphi_x + \psi) - \int_0^\infty g(s) \eta_{xx}(s)  ds - m \theta = 0 & \text{in } (0,L) \times \mathbb{R}^+, \\
		\rho_3 \theta_t + q_x + m (\varphi_x + \psi)_t = 0 & \text{in } (0,L) \times \mathbb{R}^+, \\
		\tau q_t + \sigma q + \theta_x = 0& \text{in } (0,L) \times \mathbb{R}^+,\\
		\eta_t + \eta_s - \psi_t = 0& \text{in } (0,L) \times \mathbb{R}^+,
	\end{cases}
\end{align}
for $s>0$, where the boundary satisfies the following conditions
\begin{align}\label{with memory:boundary condition}
	\begin{cases}
		u_x(0,t) = u_x(L,t) = \varphi_x(0,t) = \varphi_x(L,t) = \psi(0,t) = \psi(L,t) = \theta(0,t) = \theta(L,t) = 0, \\
		\eta^t(0,s) = \eta^t(L,s) = 0, \quad \eta^t(\cdot,0) = 0,
	\end{cases}
\end{align}
for $s, t \in \mathbb{R}^+$,
and the initial values are defined as follows
\begin{equation}\label{with memory:initial condition}
	\begin{cases}
		u(x,0) = u_0(x), \quad u_t(x,0) = v_0(x), \\
		\varphi(x,0) = \varphi_0(x), \quad \varphi_t(x,0) = \Phi_0(x), \\
		\psi(x,0) = \psi_0(x), \quad \psi_t(x,0) = \Psi_0(x), \\
		\theta(x,0) = \theta_0(x), \quad q(x,0)=q_0(x), \\
		\eta^0(x,s) = \eta_0(x,s),
	\end{cases}
\end{equation}
for $x \in (0, L)$, $s \in \mathbb{R}^+$.

We introduce the following image space
\[
\mathcal{H}_2 := H_*^1 \times L_*^2 \times H_*^1 \times L_*^2 \times H_0^1 \times L^2 \times L^2 \times L^2 \times L_g^2(\mathbb{R}^+; H_0^1),
\]
which forms a Hilbert space when endowed with the inner product
\begin{align*}
	(U_1, U_2)_{\mathcal{H}_2} = &\rho(v_1, v_2) + \alpha (u_{1x}, u_{2x}) + \rho_1(\Phi_1, \Phi_2) + \rho_2(\Psi_1, \Psi_2) + \dt b (\psi_{1x},\psi_{2x}) \\
	&+ \beta (\varphi_1 - u_1, \varphi_2 - u_2) + \kappa(\varphi_{1x} + \psi_1, \varphi_{2x} + \psi_2) + \rho_3(\theta_1, \theta_2) + \tau(q_1, q_2) + (\eta_1, \eta_2)_{L_g^2(\mathbb{R}^+; H_0^1)},
\end{align*}
and the norm
\[
\|U\|_{\mathcal{H}_2}^2 := \rho\|v\|^2 + \alpha\|u_x\|^2 + \rho_1\|\Phi\|^2 + \rho_2\|\Psi\|^2 + \dt b\|\psi_x\|^2 + \beta\|\varphi - u\|^2 + \kappa\|\varphi_x + \psi\|^2 + \rho_3\|\theta\|^2 + \tau \|q\| ^2 + \|\eta\|_{L_g^2(\mathbb{R}^+; H_0^1)} ^2,
\]
where $U := (u, v, \varphi, \Phi, \psi, \Psi, \theta, q, \eta) \in \mathcal{H}_2$ and $U_i := (u_i, v_i, \varphi_i, \Phi_i, \psi_i, \Psi_i, \theta_i, q_i, \eta_i) \in \mathcal{H}_2$ for $i = 1, 2$.
Thus, the system can be transformed into the following Cauchy problem
\begin{equation}\label{model-abstract-2}
	\begin{cases}
		\dfrac{d}{dt}U = \mathcal{A}_2U, & t > 0, \\
		U(0) =  U_0 :=(u_0, v_0, \varphi_0, \Phi_0, \psi_0, \Psi_0, \theta_0, q_0, \eta_0),
	\end{cases}
\end{equation}
where$\mathcal{A}_2: \mathcal{D}(\mathcal{A}_2) \subset \mathcal{H}_2 \to \mathcal{H}_2$ is a linear operator defined as  \begin{equation}\label{definition:A-2}
	\mathcal{A}_2 U := \begin{pmatrix}
		v \\
		\rho^{-1} \left[ \alpha u_{xx} + \beta (\varphi - u) - \gamma v \right] \\
		\Phi \\
		\rho_1^{-1} \left[ \kappa (\varphi_x + \psi)_x - \beta (\varphi - u) - m \theta_x \right] \\
		\Psi \\
		\rho_2^{-1} \left[ \dt b \psi_{xx} - \kappa (\varphi_x + \psi) + \int_0^\infty g(s) \eta_{xx}(s) ds + m\theta \right] \\
		\rho_3^{-1} \left[ -q_x - m (\Phi_x + \Psi) \right] \\
		\tau^{-1}(-\sigma q - \theta_x) \\
		\Psi - \eta_s
	\end{pmatrix}.
\end{equation}
The domain of the operator $\mathcal{A}_2$ is defined as follows
\begin{align*}
	\mathcal{D}(\mathcal{A}_2) := \bigg\{ U \in \mathcal{H}_2 :&
	v, \Phi \in H_*^1;  u_x, \varphi_x, \theta, \Psi \in H_0^1; \\
	&\dt b \psi + \int_0^\infty g(s)\eta(s)ds \in H^2; \\
	&q \in H^1;
	\eta_s \in L_g^2(\mathbb{R}^+; H_0^1) ; \eta(\cdot, 0) = 0
	\bigg\}.
\end{align*}
The main results regarding well-posedness and exponential stability are presented as follows:
\begin{theorem}[Well-posedness]\label{Well-Posedness-2}
	Let \(\mathcal{A}_2\) be the linear operator defined in \eqref{definition:A-2} with domain \(\mathcal{D}(\mathcal{A}_2) \subset \mathcal{H}_2\). Then \(\mathcal{A}_2\) serves as the infinitesimal generator of a \(C_0\)-contraction semigroup \(\{e^{t\mathcal{A}_2}\}_{t \ge 0}\) on \(\mathcal{H}_2\). Consequently:
        If \(U_0 \in \mathcal{H}_2\), the abstract Cauchy problem \eqref{model-abstract-2} admits a unique mild solution \(U(t) = e^{t\mathcal{A}_2}U_0\) satisfying \(U \in C(\mathbb{R}^+, \mathcal{H}_2)\); If \(U_0 \in \mathcal{D}(\mathcal{A}_2)\) (the domain of \(\mathcal{A}_2\)), then \(U(t) = e^{t\mathcal{A}_2}U_0\) is the unique classical solution to \eqref{model-abstract-2}, and it belongs to \(C(\mathbb{R}^+, \mathcal{D}(\mathcal{A}_2)) \cap C^1(\mathbb{R}^+, \mathcal{H}_2)\).
\end{theorem}
\begin{theorem}[Exponential stability]\label{exponential stability}
	 Under Assumption \ref{assumg}, the \(C_0\)-semigroup \(\{e^{\mathcal{A}_2 t}\}_{t \geq 0}\) generated by \(\mathcal{A}_2\) is exponentially stable. Specifically, there exist positive constants \(\alpha, \beta > 0\) (independent of \(t\)) such that
        \[
        \|e^{\mathcal{A}_2 t}\|_{\mathcal{L}(\mathcal{H}_2 \to \mathcal{H}_2)} \leq \alpha e^{-\beta t} \quad \text{for all } t \ge 0.
        \]
\end{theorem}
\subsection{Well-Posedness}\label{sec3.1}

Following the approach in Subsection \ref{sec2.1}, to get Theorem \ref{Well-Posedness-2}, we only need to show that $\mathcal{A}_2$ is both dissipative and maximal, as stated in the following two lemmas.
\begin{lemma}\label{dissipative-2}
	The operator $\mathcal{A}_2$ defined in \eqref{definition:A-2} is dissipative.
\end{lemma}
\begin{proof}
	For any $U \in \mathcal{D}(\mathcal{A}_2)$, using Assumption \ref{assumg}, a direct calculation gives
	\begin{align*}
		{\rm Re}(\mathcal{A}_2U,U)_{\mathcal{H}_2} &= {\rm Re}\bigg\{
		\rho\left(\rho^{-1} \left[ \alpha u_{xx} + \beta (\varphi - u) - \gamma v \right], v\right) + \alpha\left(v_x, u_x\right)  + \rho_1\left(\rho_1^{-1} \left[ \kappa (\varphi_x + \psi)_x - \beta (\varphi - u) - m \theta_x \right], \Phi\right) \\
		&\quad + \rho_2\left(\rho_2^{-1} \left[\dt b \psi_{xx} - \kappa (\varphi_x + \psi) + \int_0^\infty g(s) \eta_{xx}(s) ds + m\theta \right], \Psi\right) + \dt b\left(\Psi_x,\psi_x\right) + \beta\left(\Phi - v, \varphi - u\right)  \\
		&\quad + \kappa\left(\Phi_x + \Psi, \varphi_x + \psi\right)  + \rho_3\left(\rho_3^{-1} \left[ -q_x - m (\Phi_x + \Psi) \right], \theta\right)  \\
&\quad +\tau\left(\tau^{-1}(-\sigma q - \theta_x),q\right) + \int_0^\infty g(s)(\Psi_x - \eta_{sx}, \eta_x)  ds \bigg\} \\
		&= -\gamma\|v\|^2  - \sigma \|q\|^2 + \frac{1}{2}\int_{0}^{\infty}g'(s)\|\eta_x\|^2ds\\
		&\le -\gamma\|v\|^2  - \sigma \|q\|^2 - \frac{\kappa_1}{2} \int_0^\infty g(s) \|\eta_x\|^2 ds\\
		&\le 0.
	\end{align*}
	Thus, $\mathcal{A}_2$ satisfies the dissipativity condition.
\end{proof}
\begin{lemma}\label{maximal-2}
	$0 \in \varrho(\mathcal{A}_2)$ (the resolvent set of $\mathcal{A}_2$). Consequently, $\mathcal{A}_2$ is maximal.
\end{lemma}
\begin{proof}
	For $\mathcal{A}_2U = F:=(f_1,f_2,f_3,f_4,f_5,f_6,f_7,f_8,f_9) \in \mathcal{H}_2$, we aim to prove that there exists a unique solution $U\in\mathcal{D}(\mathcal{A}_2)$ satisfying $\|U\|_{\mathcal{H}_2} \leq C\|F\|_{\mathcal{H}_2}$, for some constant $C > 0$. \\
	Substitute the definition of $\mathcal{A}_2$ into $\mathcal{A}_2U = F$, we obtain the following equations
	\begin{equation}\label{eq:resolvent-system-A-2}
		\begin{cases}
			v = f_1, \\
			\alpha u_{xx} + \beta (\varphi - u) - \gamma v = \rho f_2, \\
			\Phi = f_3, \\
			\kappa (\varphi_x + \psi)_x - \beta (\varphi - u) - m \theta_x = \rho_1 f_4, \\
			\Psi = f_5, \\
			\dt b \psi_{xx} - \kappa (\varphi_x + \psi) + \int_0^\infty g(s) \eta_{xx}(s) ds + m\theta = \rho_2 f_6, \\
			-q_x - m (\Phi_x + \Psi) = \rho_3 f_7, \\
			-\sigma q - \theta_x = \tau f_8, \\
			\Psi - \eta_s = f_9.
		\end{cases}
	\end{equation}
	First, equations \eqref{eq:resolvent-system-A-2}$_{1,3,5}$ yield
	\begin{equation}\label{df1}
		v = f_1 \in H_*^1, \quad \Phi = f_3 \in H_*^1, \quad \Psi = f_5 \in H_0^1.
	\end{equation}
	Second, from equation \eqref{eq:resolvent-system-A-2}$_9$, it follows that $\eta_s = \Psi - f_9 \in L_g^2(\mathbb{R}^+; H_0^1)$. Integrating with respect to $s$ yields
	\begin{align}\label{def:eta}
		\eta(s) = s\Psi - \int_0^s f_9(\tau)d\tau\in C(\mathbb{R}^+; H_0^1).
	\end{align}
	Now, we prove that $\eta \in L_g^2(\mathbb{R}^+; H_0^1)$. First, we show, $\int_0^N g(s)\|\eta_x\|^2ds$ is bounded for any $N > 0$.
	By the expression of $\eta$ \eqref{def:eta}, the following estimate holds
	\begin{align*}
		\|\eta_x\|^2 &= \|s\Psi_x - \int_0^s f_{9x}(\tau)d\tau\|^2\\
		 &\le 2s^2\|\Psi_x\|^2 + 2s\int_0^s\| f_{9x}(\tau)\|^2d\tau  \\
		&\le 2s^2\|\Psi_x\|^2 + \frac{2s}{g(s)}\int_0^\infty g(\tau)\| f_{9x}(\tau)\|^2d\tau \\
		&= 2s^2\|\Psi_x\|^2 + \frac{2s}{g(s)}\|f_9\|_{L_g^2(\mathbb{R}^+; H_0^1)} ^2.
	\end{align*}
	Combined with the Assumption \ref{assumg}, we conclude that $g$ decays exponentially. Therefore, for any $N > 0$, there exists a constant $C_N$ depending on $N$ such that,
	\begin{align}
		\int_0^N g(s)\|\eta_x\|^2ds &\leq 2\int_0^N g(s)s^2ds \|\Psi_x\|^2 + 2\int_0^N s ds \|f_9\|_{L_g^2(\mathbb{R}^+; H_0^1)}^2\notag \\
		&\leq C_N.\label{ddf}
	\end{align}
	Now, using Assumption \ref{assumg}, integrating by parts, and applying Young's inequality, it follows that
	\begin{align*}
		\int_0^N g(s)\|\eta_x\|^2ds &\leq -\frac{1}{\kappa_1}\int_0^N g'(s)\|\eta_x\|^2ds \\
		&= -\frac{g(N)}{\kappa_1}\|\eta_x(N)\|^2 + \frac{g(0)}{\kappa_1}\|\eta_x(0)\|^2 + \frac{2}{\kappa_1}\int_0^N g(s){\rm Re}\left(\eta_x, \eta_{sx}\right)ds \\
		&\leq \frac{2}{\kappa_1}\int_0^N g(s)\|\eta_x\|\|\eta_{sx}\|ds \\
		&\leq \frac{1}{2}\int_0^N g(s)\|\eta_x\|^2ds + \frac{2}{\kappa_1^2}\int_0^N g(s)\|\eta_{sx}\|^2ds.
	\end{align*}
	Since by estimate \eqref{ddf}, $\int_0^N g(s)\|\eta_x\|^2ds$ is bounded. The terms in the above inequality can be reordered to yield
	\begin{align*}
		\int_0^N g(s)\|\eta_x\|^2ds \leq \frac{4}{\kappa_1^2}\int_0^N g(s)\|\eta_{sx}\|^2ds.
	\end{align*}
	Then, letting $N \to \infty$, we find that $\eta \in L_g^2(\mathbb{R}^+; H_0^1)$ follows directly from $\eta_s \in L_g^2(\mathbb{R}^+; H_0^1)$.
	
	Using equation \eqref{eq:resolvent-system-A-2}$_8$, substituting $q$ into equation \eqref{eq:resolvent-system-A-2}$_7$, and rearranging equations \eqref{eq:resolvent-system-A-2}$_{2,4,6}$, we derive
	\begin{align}\label{eq:resolvent-system-1-A-2}
		\begin{cases}
			-\alpha u_{xx} - \beta(\varphi - u) = -\rho f_2 - \gamma v=-\rho f_2 - \gamma f_1 := g_1 \in L_*^2, \\
			-\kappa(\varphi_x + \psi)_x + \beta(\varphi - u) + m\theta_x = -\rho_1 f_4 := g_2 \in L_*^2, \\
			- \dt b \psi_{xx} + \kappa(\varphi_x + \psi) - m\theta = -\rho_2 f_6 +  \int_0^\infty g(s) \eta_{xx}(s) ds  := g_3 \in H^{-1}, \\
			- \theta_{xx} = - \sigma m(\Phi_x + \Psi) - \sigma \rho_3 f_7 + \tau f_{8x}=- \sigma m(f_{3x} + f_5) - \sigma \rho_3 f_7 + \tau f_{8x}:=g_4 \in H^{-1}.
		\end{cases}
	\end{align}
	Following the method of Lemma \ref{maximal-1}, one can prove that equations \eqref{eq:resolvent-system-1-A-2}  has a unique solution $(u,\varphi,\psi,\theta) \in H_*^1 \times H_*^1 \times H_0^1 \times H_0^1$. Furthermore, we can solve for
	$q$, demonstrate that $U \in \mathcal{D}(\mathcal{A}_2)$ and $\|U\|_{\mathcal{H}_2} \leq C\|F\|_{\mathcal{H}_2}$ similarly. Consequently, $0 \in \varrho(\mathcal{A}_2)$.
\end{proof}

\subsection{Exponential Decay} \label{sec3.2}
In this section, we prove the exponential stability result, i.e., Theorem \ref{exponential stability}. The proof is carried out in two sequential steps, as detailed below. Here we invoke the following lemma

\begin{lemma}\label{Thm:SpectralApprox}\cite{MR933321}
	Let $\mathcal{A}$ be the infinitesimal generator of a $C_0$-semigroup on a reflexive Banach space $X$. Then
	$$
	\sigma(\mathcal{A}) \cap i\mathbb{R} = \sigma_{ap}(\mathcal{A}) \cap i\mathbb{R},
	$$
	where $\sigma(\mathcal{A})$ and $\sigma_{ap}(\mathcal{A})$ denote the spectrum and approximate spectrum of the operator $\mathcal{A}$, respectively, and $\sigma_{ap}(\mathcal{A})$ is defined as follows
\begin{align}\label{defapp}
	\sigma_{ap}(\mathcal{A}):= \{ \lambda \in \mathbb{C}: \exists U_{n} \in \mathcal{D}(\mathcal{A}), \|U_{n}\|_{\mathcal{H}} = 1 \text{ and } (\lambda I - \mathcal{A})U_{n} \to 0 \text{ in } \mathcal{H}\},
\end{align}
\end{lemma}

\textbf{Step1: The Resolvent Set of $\mathcal{A}_2$ Contains the Imaginary Axis.}
To prove that $i\mathbb{R} \subset \rho(\mathcal{A}_2)$, we employ a contradiction argument. Suppose, for contradiction, that $i\mathbb{R} \not\subset \rho(\mathcal{A}_2)$. Since $0 \in \rho(\mathcal{A}_2)$(Lemma \ref{maximal-2}), we assume without loss of generality that there exists $\omega > 0$ such that $i\omega \notin \rho(\mathcal{A}_2)$. Thus, $i\omega \in \sigma(\mathcal{A}_2)$. Then by Lemma \ref{Thm:SpectralApprox}, $i\omega \in \sigma_{ap}(\mathcal{A}_2)$. From the definition of $\sigma_{ap}(\mathcal{A}_2)$, it follows that
there exists a sequence $\{U_n=(u_n, v_n, \varphi_n, \Phi_n, \psi_n, \Psi_n, \theta_n, q_n, \eta_n) \} \subset \mathcal{D}(\mathcal{A}_2)$ with $\|U_n\|_{\mathcal{H}_2} = 1$ such that
\begin{align}\label{convergence}
	(i\omega I - \mathcal{A}_2)U_n \to 0 ~~~~~~~~~\text{in} ~\mathcal{H}_2.
\end{align}
Using the definition of $\mathcal{A}_2$ \eqref{definition:A-2}, the convergence \eqref{convergence} yields
\begin{align}\label{omega-limit-system}
	\begin{cases}
		i\omega u_n - v_n  \to 0 \quad &\text{in } H_*^1,\\
		i\omega \rho v_n -\alpha u_{nxx}-\beta(\varphi_n - u_n)+ \gamma v_n \to 0 \quad &\text{in } L_*^2,\\
		i\omega\varphi_n - \Phi_n \to 0\quad &\text{in } H_*^1,\\
		i\omega\rho_1 \Phi_n - \kappa(\varphi_{nx}+\psi_n)_x + \beta (\varphi_n - u_n) +  m\theta_{nx} \to 0 \quad &\text{in } L_*^2,\\
		i\omega \psi_n - \Psi_n \to 0 \quad &\text{in } H_0^1,\\
		i\omega \rho_2 \Psi_n - \dt b \psi_{nxx} + \kappa(\varphi_{nx} + \psi_n) - \int_{0}^{\infty} g(s) \eta_{nxx}(s) ds - m\theta_n \to 0 \quad &\text{in } L^2,\\
		i\omega \rho_3 \theta_n + q_{nx} + m(\Phi_{nx} + \Psi_n) \to 0 \quad &\text{in } L^2,\\
		i\omega\tau q_n + \sigma q_n + \theta_{nx} \to 0 \quad &\text{in } L^2,\\
		i\omega \eta_n - \Psi_n + \eta_{ns} \to 0\quad &\text{in } L_g^2(\mathbb{R}^+;H_0^1).
	\end{cases}
\end{align}
\begin{lemma}\label{iR:vn,qn,etan}
	Under Assumption \ref{assumg}, as $n \to \infty$, we have
	\begin{align}
		\|v_n\| \to 0,\label{limit:vn}\\
		\|q_n\| \to 0,\label{limit:qn}\\
		\int_{0}^{\infty} g'(s)\|\eta_{nx}(s)\|^2ds \to 0,\label{limit:g'}\\
		\int_{0}^{\infty} g(s)\|\eta_{nx}(s)\|^2ds \to 0.\label{limit:g}
	\end{align}
\end{lemma}

\begin{proof}
	Take the real part of the $\mathcal{H}_2$-inner product of $i \omega U_n - \mathcal{A}_2 U_n$ with $U_n$
	\begin{align*}
		\operatorname{Re} ( i \omega U_n - \mathcal{A}_2 U_n, U_n )_{\mathcal{H}_2}
		&= -\operatorname{Re} ( \mathcal{A}_2 U_n, U_n )_{\mathcal{H}_2}\notag\\
		&= \gamma\|v_n\|^2 + \sigma\|q_n\|^2 - \frac{1}{2}\int_{0}^{\infty} g'(s)\|\eta_{nx}\|^2ds.
	\end{align*}
	From the convergence relation \eqref{convergence} and $\|U_n\|_{\mathcal{H}_2} = 1$, we have
	\begin{align*}
		\operatorname{Re} ( i \omega U_n - \mathcal{A}_2 U_n, U_n )_{\mathcal{H}_2} \to 0.
	\end{align*}
	Thus, \eqref{limit:vn}--\eqref{limit:g'} follow immediately. Estimate \eqref{limit:g} is derived from \eqref{limit:g'} using Assumption \ref{assumg}.
\end{proof}
\begin{lemma}\label{iR:thetan}
	Under Assumption \ref{assumg}, as $n \to \infty$, we have
	\begin{align}
		\|\theta_{nx}\| \to 0.\label{limit:thetan}
	\end{align}
\end{lemma}
\begin{proof}
	From \eqref{omega-limit-system}$_8$ and \eqref{limit:qn}, it follows that
	\begin{align*}
		\|\theta_{nx}\| \leq \|i\omega\tau q_n + \sigma q_n + \theta_{nx}\| + \|i\omega\tau q_n + \sigma q_n\| \to 0.
	\end{align*}
\end{proof}
\begin{lemma}\label{iR:un}
	Under Assumption \ref{assumg}, as $n \to \infty$, we have
	\begin{align}
		\|u_{nx}\| \to 0.\label{limit:un}
	\end{align}
\end{lemma}
\begin{proof}
	By \eqref{omega-limit-system}$_1$ and \eqref{limit:vn}, we obtain
	\begin{align*}
		\omega\|u_n\| = \|i\omega u_n\| \leq \|i\omega u_n - v_n\| + \|v_n\| \to 0.
	\end{align*}
	Thus $\|u_n\| \to 0$.
	Take the $L^2$-inner product of \eqref{omega-limit-system}$_2$ with $u_n$ and integrate by parts
	\begin{align*}
		(i\omega\rho v_n - \beta(\varphi_n - u_n) + \gamma v_n, u_n) + \alpha\|u_{nx}\|^2 \rightarrow 0.
	\end{align*}
	The first term satisfies the following estimate
	\begin{align*}
		|(i\omega\rho v_n - \beta(\varphi_n - u_n) + \gamma v_n, u_n)| \leq (( \omega\rho + \gamma ) \|v_n\| + \beta \|\varphi_n - u_n\|) \|u_n\| \leq c \|u_n\| \rightarrow 0.
	\end{align*}
	Therefore, $\|u_{nx}\| \to 0$.
\end{proof}
\begin{lemma}\label{iR:Phin}
	Under Assumption \ref{assumg}, as $n \to \infty$, we have
	\begin{align}
		\|\Phi_n\| \to 0.\label{limit:Phin}
	\end{align}
\end{lemma}
\begin{proof}
	Take the $L^2$-inner product of \eqref{omega-limit-system}$_2$ with $\varphi_n$ and integrate by parts
	\begin{align}\label{jvarphin}
		(i \omega \rho v_n + \beta u_n + \gamma v_n, \varphi_n) + \alpha(u_{nx}, \varphi_{nx}) - \beta \|\varphi_n\|^2 \to 0.
	\end{align}
	The sum of the first two terms satisfies the following estimate
	\begin{align*}
		&|(i \omega \rho v_n + \beta u_n + \gamma v_n, \varphi_n) + \alpha(u_{nx}, \varphi_{nx})| \\
		&\leq (( \omega \rho + \gamma) \|v_n\| + \beta \|u_n\|) \|\varphi_n\| + \alpha \|u_{nx}\| \|\varphi_{nx}\| \\
		&\leq C(\|v_n\| + \|u_n\| + \|u_{nx}\|) \\
		&\to 0.
	\end{align*}
	Returning to \eqref{jvarphin}, we obtain: $\|\varphi_n\| \to 0$.
	Combined with \eqref{omega-limit-system}$_3$, it follows that $\|\Phi_n\| \to 0$.
\end{proof}
\begin{lemma}\label{iR:varphin}
	Under Assumption \ref{assumg}, as $n \to \infty$, we have
	\begin{align}
		\|\varphi_{nx} + \psi_n\| \to 0.\label{limit:varphin}
	\end{align}
\end{lemma}
\begin{proof}
	Using $\|\Phi_n\|, \|\varphi_n\|, \|u_n\|, \|\theta_{nx}\| \to 0$, from \eqref{omega-limit-system}$_4$, we get
	\begin{align*}
		\|(\varphi_{nx} + \psi_n)_x\| \to 0.
	\end{align*}
	Since $\varphi_{nx} + \psi_n \in H_0^1$, Poincar\'e's inequality yields
	\begin{align*}
		\|\varphi_{nx} + \psi_n\| \leq \sqrt{c_p}\|(\varphi_{nx} + \psi_n)_x\| \to 0.
	\end{align*}
\end{proof}
\begin{lemma}\label{iR:psin,Psin}
	Under Assumption \ref{assumg}, as $n \to \infty$, we have
	\begin{align}
		\|\psi_{nx}\| \to 0,\label{limit:psin}\\
		\|\Psi_{nx}\| \to 0.\label{limit:Psin}
	\end{align}
\end{lemma}
\begin{proof}
	Take the $L^2$-inner product of \eqref{omega-limit-system}$_6$ with $\psi_n$ and integrate by parts
	\begin{align}\label{eq:psin}
		i\omega\rho_2(\Psi_n, \psi_n) + \dt b\|\psi_{nx}\|^2 + (\kappa(\varphi_{nx} + \psi_n) - m\theta_n, \psi_n) + \left(\int_{0}^{\infty} g(s)\eta_{nx}(s)ds,  \psi_{nx}\right) \to 0.
	\end{align}
	First, estimate the term $(\kappa(\varphi_{nx} + \psi_n) - m\theta_n, \psi_n)$
	\begin{align*}
		\left|(\kappa(\varphi_{nx} + \psi_n) - m\theta_n, \psi_n)\right| \leq (\kappa \|\varphi_{nx} + \psi_n\|  + |m|\|\theta_n\|)\|\psi_{n}\|
		\leq C(\|\varphi_{nx} + \psi_n\| + \|\theta_n\|)
		\to 0.
	\end{align*}
	Second, estimate the term $\left(\int_{0}^{\infty} g(s)\eta_{nx}(s)ds,  \psi_{nx}\right)$
	\begin{align*}
		\left|\left(\int_{0}^{\infty} g(s)\eta_{nx}(s)ds,  \psi_{nx}\right)\right| \leq \left\|\int_{0}^{\infty} g(s)\eta_{nx}(s)ds\right\| \|\psi_{nx}\|
		\leq b_0^{\frac{1}{2}} \|\eta_n\|_{L^2_g(\mathbb{R}^+; H_0^1)} \|\psi_{nx}\|
		\leq C \|\eta_n\|_{L^2_g(\mathbb{R}^+; H_0^1)}
		\to 0.
	\end{align*}
	Third, estimate the term $(\Psi_n, \psi_n)$. Taking the $L^2$-inner product of \eqref{omega-limit-system}$_7$ with $\psi_n$, we obtain
	\begin{align}\label{eq:Psin,psin}
		i\omega\rho_3(\theta_n, \psi_n) - (q_n, \psi_{nx}) - m(\Phi_n, \psi_{nx}) + m(\Psi_n, \psi_n) \to 0.
	\end{align}
	From \eqref{eq:Psin,psin}, \eqref{limit:thetan}, \eqref{limit:qn}, and \eqref{limit:Phin}, we derive
	\begin{align*}
		|m(\Psi_n, \psi_n)| &\leq |i\omega\rho_3(\theta_n, \psi_n) - (q_n, \psi_{nx}) - m(\Phi_n, \psi_{nx}) + m(\Psi_n, \psi_n)|  + |i\omega\rho_3(\theta_n, \psi_n) - (q_n, \psi_{nx}) - m(\Phi_n, \psi_{nx})| \\
		&\leq |i\omega\rho_3(\theta_n, \psi_n) - (q_n, \psi_{nx}) - m(\Phi_n, \psi_{nx}) + m(\Psi_n, \psi_n)|  + \omega\rho_3 \|\theta_n\|\|\psi_n\| + \|q_n\|\|\psi_{nx}\| + |m|\|\Phi_n\|\|\psi_{nx}\| \\
		&\to 0.
	\end{align*}
	Returning to \eqref{eq:psin} yields $\|\psi_{nx}\| \to 0$. $\|\Psi_{nx}\| \to 0$ is a direct consequence of \eqref{omega-limit-system}$_5$ and \eqref{limit:psin}.
\end{proof}

From Lemmas \ref{iR:vn,qn,etan}-\ref{iR:psin,Psin}, it follows that $\|U_n\|_{\mathcal{H}_2} \to 0$, contradicting $\|U_n\|_{\mathcal{H}_2} = 1$. Thus, we derive $i \mathbb{R} \subset \varrho(\mathcal{A}_2)$.

\textbf{Step2: Uniform Boundedness of $(i\lambda I-\mathcal{A}_2)^{-1}$.} Consider the following resolvent equation
\begin{align}\label{resolvent eq}
	(i\lambda I-\mathcal{A}_2)U=F := (f_1,f_2,f_3,f_4,f_5,f_6,f_7,f_8,f_9) \in \mathcal{H}_2,
\end{align}
where $U=(u,v,\varphi,\Phi,\psi,\Psi,\theta,q,\eta) \in \mathcal{D}(\mathcal{A}_2)$. From the definition of $\mathcal{A}_2$ \eqref{definition:A-2}, equation \eqref{resolvent eq} gives the following system

\begin{align}\label{Resolvent equations-2}
	\begin{cases}
		i\lambda u - v=f_1   \quad &\text{in } H_*^1,\\
		i\lambda \rho v -\alpha u_{xx}-\beta(\varphi - u)+ \gamma v =\rho f_2 \quad &\text{in } L_*^2,\\
		i\lambda\varphi - \Phi =f_3\quad &\text{in } H_*^1,\\
		i\lambda\rho_1 \Phi - \kappa(\varphi_{x}+\psi)_x + \beta (\varphi - u) +  m\theta_{x}=\rho_1f_4  \quad &\text{in } L_*^2,\\
		i\lambda \psi - \Psi =f_5 \quad &\text{in } H_0^1,\\
		i\lambda \rho_2 \Psi - \dt b \psi_{xx} + \kappa(\varphi_{x} + \psi) - \int_{0}^{\infty} g(s) \eta_{xx}(s) ds - m\theta =\rho_2f_6 \quad &\text{in } L^2,\\
		i\lambda \rho_3 \theta + q_x + m(\Phi_{x} + \Psi) =\rho_3f_7 \quad &\text{in } L^2,\\
		i\lambda \tau q + \sigma q + \theta_x = \tau f_8 \quad &\text{in } L^2, \\
		i\lambda \eta + \eta_s - \Psi=f_9 \quad &\text{in } L_g^2(\mathbb{R}^+;H_0^1).
	\end{cases}
\end{align}
\begin{lemma}\label{boundedness:v,q,eta}
	Under Assumption \ref{assumg}, the following estimates hold
	\begin{align}
		\|v\|^2 \leq \frac{1}{\gamma}\|U\|_{\mathcal{H}_2}\|F\|_{\mathcal{H}_2}, \label{v}\\
		\|q\|^2 \leq \frac{1}{\sigma}\|U\|_{\mathcal{H}_2}\|F\|_{\mathcal{H}_2}, \label{q}\\
		-\int_{0}^{\infty}g'(s)\|\eta_x(s)\|^2ds \leq 2\|U\|_{\mathcal{H}_2}\|F\|_{\mathcal{H}_2},\label{g'}\\
		\|\eta\|_{L_g^2(\mathbb{R}^+; H_0^1) }^2 \leq \frac{2}{\kappa_1}\|U\|_{\mathcal{H}_2}\|F\|_{\mathcal{H}_2}.\label{g}
	\end{align}
\end{lemma}
\begin{proof}
	Taking the $\mathcal{H}_2$-inner product of $(i\lambda I-\mathcal{A}_2)U=F$ with $U$, we derive
	\begin{align*}
		\|F\|_{\mathcal{H}_2}\|U\|_{\mathcal{H}_2}
		&\geq |\operatorname{Re}(F,U)_{\mathcal{H}_2}| = |\operatorname{Re}(i\lambda U-\mathcal{A}_2U,U)_{\mathcal{H}_2}| \\
		& = |-\operatorname{Re}(\mathcal{A}_2U,U)_{\mathcal{H}_2}| = \gamma \|v\|^2 + \sigma \|q\|^2 - \frac{1}{2}\int_{0}^{\infty} g'(s)\|\eta_x(s)\|^2ds.
	\end{align*}
	Estimates \eqref{v}--\eqref{g'} follow immediately. Estimate \eqref{g} is derived from \eqref{g'} using Assumption \ref{assumg}.
\end{proof}

Following the same procedure as in the proofs of Lemmas \ref{lemma-ux}-\ref{lemma-Phi}, we obtain the following lemma\\
\begin{lemma}\label{boundedness:ux}
	Under Assumption \ref{assumg}, for any $\epsilon>0$, there exists a constant $C_{\epsilon}>0$ (independent of $\lambda$) such that
	\begin{align}
		\|u_x\|^2, \|\theta\|^2, \|\varphi_x + \psi\|^2, \|\varphi - u\|^2, \|\Phi\|^2 \leq \epsilon\|U\|_{\mathcal{H}_2}^2 + C_{\epsilon}\|F\|_{\mathcal{H}_2}^2
	\end{align}
	for $|\lambda| \geq 1$.
\end{lemma}

\begin{lemma}\label{boundedness:Psi}
	Under Assumption \ref{assumg}, for any $\epsilon>0$, there exists a constant $C_{\epsilon}>0$ (independent of $\lambda$) such that
	\begin{align}
		 \|\Psi\|^2 \leq \epsilon\|U\|_{\mathcal{H}_2}^2 + C_{\epsilon}\|F\|_{\mathcal{H}_2}^2
	\end{align}
	for $|\lambda| \geq 1$.
\end{lemma}
\begin{proof}
	Take the $L^2$-inner product of \eqref{Resolvent equations-2}$_6$ with $\Psi$ and integrate by parts
	\begin{align}\label{hh}
		i\lambda\rho_2 \|\Psi\|^2 + \dt b(\psi_x, \Psi_x) + \kappa(\varphi_x + \psi, \Psi) + \int_{0}^{\infty} g(s)(\eta_x, \Psi_x) ds - m(\theta, \Psi) = \rho_2 (f_6, \Psi).
	\end{align}
	Multiplying both sides of the above equation by $ g(s) $ and integrating with respect to $ s $ from $(0,\infty)$, we obtain
	\begin{align}\label{eq:Psi}
		 i\lambda\rho_2 b_0 \|\Psi\|^2 + \dt b(\psi, \Psi)_{L_g^2(\mathbb{R}^+; H_0^1)} + \kappa b_0(\varphi_x + \psi, \Psi) + b_0 (\eta, \Psi)_{L_g^2(\mathbb{R}^+; H_0^1)} - m b_0 (\theta, \Psi) = \rho_2 b_0(f_6, \Psi).
	\end{align}
	First, we estimate the term $	(\psi, \Psi)_{L_g^2(\mathbb{R}^+; H_0^1)}$. Using \eqref{Resolvent equations-2}$_9$, we obtain
	\begin{align*}
		(\psi, \Psi)_{L_g^2(\mathbb{R}^+; H_0^1)} &= (\psi, i\lambda \eta + \eta_s - f_9)_{L_g^2(\mathbb{R}^+; H_0^1)} \\
		&= -i\lambda (\psi, \eta)_{L_g^2(\mathbb{R}^+; H_0^1)} + (\psi, \eta_s)_{L_g^2(\mathbb{R}^+; H_0^1)} - (\psi, f_9)_{L_g^2(\mathbb{R}^+; H_0^1)} \\
		&= -i\lambda (\psi, \eta)_{L_g^2(\mathbb{R}^+; H_0^1)} - \int_0^\infty g'(s)(\psi_x, \eta_x) ds - (\psi, f_9)_{L_g^2(\mathbb{R}^+; H_0^1)}.
	\end{align*}
	By H\"older's and Young's inequalities, using the estimates \eqref{g} and
	\eqref{g'} yields
	\begin{align*}
		|(\psi, \Psi)_{L_g^2(\mathbb{R}^+; H_0^1)}| &\leq |\lambda| \|\psi\|_{L_g^2(\mathbb{R}^+; H_0^1)}\|\eta\|_{L_g^2(\mathbb{R}^+; H_0^1)} - \int_0^\infty g'(s) \|\eta_x\| ds \|\psi_x\| + \|\psi\|_{L_g^2(\mathbb{R}^+; H_0^1)} \|f_9\|_{L_g^2(\mathbb{R}^+; H_0^1)} \\
		&\leq |\lambda| b_0^{\frac{1}{2}} \|\psi_x\| \|\eta\|_{L_g^2(\mathbb{R}^+; H_0^1)} + g(0)^{\frac{1}{2}} \left( \int_0^\infty (-g'(s)) \|\eta_x\|^2 ds \right) ^{\frac{1}{2}} \|\psi_x\| +  b_0^{\frac{1}{2}}\|\psi_x\|\|f_9\|_{L_g^2(\mathbb{R}^+; H_0^1)}\\
		&\leq |\lambda| \left( \epsilon \|U\|_{\mathcal{H}_2}^2 + C_{\epsilon} \|F\|_{\mathcal{H}_2}^2 \right).
	\end{align*}
	Second, we estimate the term $(\eta, \Psi)_{L_g^2(\mathbb{R}^+; H_0^1)}$. From equation \eqref{Resolvent equations-2}$_5$, it follows that
	\begin{align*}
		|(\eta, \Psi)_{L_g^2(\mathbb{R}^+; H_0^1)}| &= |(\eta, i\lambda\psi - f_5)_{L_g^2(\mathbb{R}^+; H_0^1)}| \\
		&\leq \|\eta\|_{L_g^2(\mathbb{R}^+; H_0^1)} \|i\lambda\psi - f_5\|_{L_g^2(\mathbb{R}^+; H_0^1)} \\
		&\leq \|\eta\|_{L_g^2(\mathbb{R}^+; H_0^1)} (|\lambda| \|\psi\|_{L_g^2(\mathbb{R}^+; H_0^1)} + \|f_5\|_{L_g^2(\mathbb{R}^+; H_0^1)}) \\
		&= b_0^{\frac{1}{2}} |\lambda| \|\eta\|_{L_g^2(\mathbb{R}^+; H_0^1)} \|\psi_x\| + b_0^{\frac{1}{2}} \|\eta\|_{L_g^2(\mathbb{R}^+; H_0^1)} \|f_{5x}\| \\
		&\leq |\lambda| \left( \epsilon \|U\|_{\mathcal{H}_2}^2 + C_{\epsilon} \|F\|_{\mathcal{H}_2}^2 \right).
	\end{align*}
	Returning to equation \eqref{eq:Psi}, we get $\|\Psi\|^2 \leq \epsilon\|U\|_{\mathcal{H}_2}^2 + C_{\epsilon}\|F\|_{\mathcal{H}_2}^2$, for $|\lambda| \geq 1$ immediately.
\end{proof}

\begin{lemma}\label{boundedness:psix}
	Under Assumption \ref{assumg}, for any $\epsilon>0$, there exists a constant $C_{\epsilon}>0$ (independent of $\lambda$) such that
	\begin{align}
		\|\psi_x\|^2 \leq \epsilon\|U\|_{\mathcal{H}_2}^2 + C_{\epsilon}\|F\|_{\mathcal{H}_2}^2
	\end{align}
	for $|\lambda| \geq 1$.
\end{lemma}
\begin{proof}
	From equation \eqref{Resolvent equations-2}$_5$, we get $\Psi_x = i\lambda \psi_x - f_{5x}$. Substituting $\Psi_x$ into the term $\dt b (\psi_x, \Psi_x)$ of equation \eqref{hh}, it follows that
	\begin{align}\label{eqpsix}
		i\lambda\rho_2 \|\Psi\|^2 - i\lambda \dt b \|\psi_x\|^2 - \dt b (\psi_x,f_{5x})+ \kappa(\varphi_x + \psi, \Psi) + (\eta,\Psi)_{L_g^2(\mathbb{R}^+;H_0^1)} - m(\theta, \Psi) = \rho_2 (f_6, \Psi).
	\end{align}
	From the proof of Lemma \ref{boundedness:Psi}, we obtain  $|(\eta, \Psi)_{L_g^2(\mathbb{R}^+; H_0^1)}| \leq |\lambda| \left( \epsilon \|U\|_{\mathcal{H}_2}^2 + C_{\epsilon} \|F\|_{\mathcal{H}_2}^2 \right).$ Then combined with Lemmas \ref{boundedness:ux} and \ref{boundedness:Psi}, from equation \eqref{eqpsix}, it follows that
	\begin{align*}
		\|\psi_x\|^2 \leq \epsilon\|U\|_{\mathcal{H}_2}^2 + C_{\epsilon}\|F\|_{\mathcal{H}_2}^2
	\end{align*}
	for $|\lambda| \geq 1$.
\end{proof}

Consequently, from Lemmas \ref{boundedness:v,q,eta}-\ref{boundedness:psix}, we derive $\|U\|_{\mathcal{H}_2}=\|(i\lambda I-\mathcal{A}_2)^{-1}F\|_{\mathcal{H}_2} \leq C\|F\|_{\mathcal{H}_2}$. The proof of this section is completed.

\begin{proof}[Proof of Theorem \ref{exponential stability}]
	By Theorem \ref{Stability Theorem}, Theorem \ref{exponential stability} is a direct consequence of steps 1 and 2.
\end{proof}

%\bibliographystyle{plain}
%\bibliography{Reference-Suspension-Bridge-Cattaneo-Heat}
\begin{thebibliography}{10}

\bibitem{MR4342054}
Mounir Afilal, Baowei Feng, and Abdelaziz Soufyane.
\newblock Optimal decay rates of a nonlinear suspension bridge with memories.
\newblock {\em Math. Methods Appl. Sci.}, 44(17):13170--13185, 2021.

\bibitem{MR4891851}
Mohammad~M. Al-Gharabli, Adel~M. Al-Mahdi, and Ahmad Mugbil.
\newblock Effects of viscoelastic damping and nonlinear feedback modulated by
  time-dependent coefficient in a suspension bridge system: existence and
  stability decay results.
\newblock {\em Mediterr. J. Math.}, 22(3):Paper No. 70, 27, 2025.

\bibitem{MR4694563}
Adel~M. Al-Mahdi, Mohammad~M. Al-Gharabli, and Mostafa Zahri.
\newblock Theoretical and numerical decay results of a viscoelastic suspension
  bridge with variable exponents nonlinearity.
\newblock {\em Math. Nachr.}, 296(12):5426--5453, 2023.

\bibitem{MR933321}
W.~Arendt and C.~J.~K. Batty.
\newblock Tauberian theorems and stability of one-parameter semigroups.
\newblock {\em Trans. Amer. Math. Soc.}, 306(2):837--852, 1988.

\bibitem{MR3023261}
I.~Bochicchio, C.~Giorgi, and E.~Vuk.
\newblock Asymptotic dynamics of nonlinear coupled suspension bridge equations.
\newblock {\em J. Math. Anal. Appl.}, 402(1):319--333, 2013.

\bibitem{MR4149715}
Ivana Bochicchio, Marco Campo, Jos\'e{}~R. Fern\'andez, and Maria~Grazia Naso.
\newblock Analysis of a thermoelastic {T}imoshenko beam model.
\newblock {\em Acta Mech.}, 231(10):4111--4127, 2020.

\bibitem{MR3240285}
Ivana Bochicchio, Claudio Giorgi, and Elena Vuk.
\newblock Long-term dynamics of a viscoelastic suspension bridge.
\newblock {\em Meccanica}, 49(9):2139--2151, 2014.

\bibitem{MR2606945}
Alexander Borichev and Yuri Tomilov.
\newblock Optimal polynomial decay of functions and operator semigroups.
\newblock {\em Math. Ann.}, 347(2):455--478, 2010.

\bibitem{MR4848548}
Meriem Chabekh, Nadhir Chougui, and Delfim F.~M. Torres.
\newblock Analysis of a shear beam model with suspenders in thermoelasticity of
  type {III}.
\newblock {\em J. Comput. Appl. Math.}, 463:Paper No. 116471, 22, 2025.

\bibitem{MR4252268}
Zayd Hajjej.
\newblock General decay of solutions for a viscoelastic suspension bridge with
  nonlinear damping and a source term.
\newblock {\em Z. Angew. Math. Phys.}, 72(3):Paper No. 90, 26, 2021.

\bibitem{MR3575770}
Jum-Ran Kang.
\newblock Asymptotic behavior of the thermoelastic suspension bridge equation
  with linear memory.
\newblock {\em Bound. Value Probl.}, pages Paper No. 206, 18, 2016.

\bibitem{MR3458828}
Jum-Ran Kang.
\newblock Global attractor for suspension bridge equations with memory.
\newblock {\em Math. Methods Appl. Sci.}, 39(4):762--775, 2016.

\bibitem{MR1681343}
Zhuangyi Liu and Songmu Zheng.
\newblock {\em Semigroups associated with dissipative systems}, volume 398 of
  {\em Chapman \& Hall/CRC Research Notes in Mathematics}.
\newblock Chapman \& Hall/CRC, Boca Raton, FL, 1999.

\bibitem{MR2142424}
Qiaozhen Ma and Chengkui Zhong.
\newblock Existence of global attractors for the coupled system of suspension
  bridge equations.
\newblock {\em J. Math. Anal. Appl.}, 308(1):365--379, 2005.

\bibitem{MR2514725}
Qiaozhen Ma and Chengkui Zhong.
\newblock Existence of strong solutions and global attractors for the coupled
  suspension bridge equations.
\newblock {\em J. Differential Equations}, 246(10):3755--3775, 2009.

\bibitem{MR3609015}
Salim~A. Messaoudi, Ahmed Bonfoh, Soh~E. Mukiawa, and Cyril~D. Enyi.
\newblock The global attractor for a suspension bridge with memory and
  partially hinged boundary conditions.
\newblock {\em ZAMM Z. Angew. Math. Mech.}, 97(2):159--172, 2017.

\bibitem{MR4887867}
L.~G.~R. Miranda, A.~J.~A. Ramos, C.~A. Nonato, C.~A. Raposo, and F.~A.~F.
  Nascimento.
\newblock Suspension bridges with thermoelasticity acting on the shear force.
\newblock {\em Z. Angew. Math. Phys.}, 76(3):Paper No. 84, 18, 2025.

\bibitem{MR4875485}
L.~G.~R. Miranda, C.~A. Raposo, and M.~M. Freitas.
\newblock Global and exponential attractors for a suspension bridge model with
  nonlinear damping.
\newblock {\em J. Differential Equations}, 431:Paper No. 113217, 20, 2025.

\bibitem{MR3754568}
Soh~Edwin Mukiawa.
\newblock Asymptotic behaviour of a suspension bridge problem.
\newblock {\em Arab J. Math. Sci.}, 24(1):31--42, 2018.

\bibitem{MR4913141}
Soh~Edwin Mukiawa, Johnson~D. Audu, and Salim~A. Messaoudi.
\newblock Stabilization of a coupled bridge system with past history and
  {G}urtin-{P}ipkin's heat conduction.
\newblock {\em Quaest. Math.}, 48(5):729--746, 2025.

\bibitem{MR4539428}
Soh~Edwin Mukiawa, Cyril~Dennis Enyi, and Salim~A. Messaoudi.
\newblock Stability of thermoelastic {T}imoshenko beam with suspenders and
  time-varying feedback.
\newblock {\em Adv. Contin. Discrete Models}, pages Paper No. 7, 19, 2023.

\bibitem{MR4886591}
Soh~Edwin Mukiawa and Salim~A. Messaoudi.
\newblock On the well-posedness and long time dynamics for a coupled nonlinear
  bridge system with past history.
\newblock {\em Appl. Math. Optim.}, 91(3):Paper No. 53, 29, 2025.

\bibitem{MR710486}
A.~Pazy.
\newblock {\em Semigroups of linear operators and applications to partial
  differential equations}, volume~44 of {\em Applied Mathematical Sciences}.
\newblock Springer-Verlag, New York, 1983.

\bibitem{MR4688212}
Carlos Raposo, Leandro Correia, Joilson Ribeiro, and Arthur Cunha.
\newblock Suspension bridge with internal damping.
\newblock {\em Acta Mech.}, 235(1):203--214, 2024.

\bibitem{timoshenko1921lxvi}
Stephen~P Timoshenko.
\newblock On the correction for shear of the differential equation for
  transverse vibrations of prismatic bars.
\newblock {\em The London, Edinburgh, and Dublin Philosophical Magazine and
  Journal of Science}, 41(245):744--746, 1921.

\end{thebibliography}

\end{document}

